\appendix
\section{Credited proofs of parameter equivalence and deformation results}
\subsection{The canonical divisor and the exact anticanonical ring}
\label{wb:sec:canonical-ring-recovery}

The analytic construction admits a further explicit calculation which does
not use its claimed identification with a sphere. Retain the actual proper
map $f:X\to B=\mathbb P^1$, the cusp $p_0$, the finite points $p_1,p_2$,
the reduced finite fibres $S_1,S_2$, and
$f^*(p_1)=3S_1$, $f^*(p_2)=4S_2$. All line-bundle isomorphisms below
refer to this constructed analytic space. No replacement period matrix or
new compactification is introduced.

\begin{theorem}[Canonical bundle and recovery of the original fibration]
\label{wb:thm:canonical-ring-recovery}
\StatementRecord{wb:thm:canonical-ring-recovery}{proved}{Historical events: GPT-6 Astra (OpenAI), ChatGPT Pro mode, effort Pro, S6\_09\_PAPERS, 17 November 2026: targeted proof check. Also checked in this writing. Pre-continuation writing events: GPT-6 Astra (OpenAI), ChatGPT Pro mode, effort Pro, S6\_10\_PAPERS, 18 November 2026: targeted proof check. Continuation metadata events: GPT-6 Astra (OpenAI), ChatGPT Pro mode, effort Pro, S6\_10\_PAPERS continuation, 19 November 2026: used and rechecked in this run; status-vocabulary reconciliation; metadata reconciliation only; no new mathematical check. \emph{Pre-continuation recorded proof scope.} Re-derived the full receiving proof: invariant local evaluation exponents0 and2; actual meromorphic pushforward floor(k/4); every graded basis, the weighted polynomial ring and the intrinsic degree-two pencil, including dependence on scalar trivializations.}{the hypotheses, notation and definitions stated here; \cref{wb:lem:hodge-line-audit}}

The evaluation map
\[
 f^*(f_*\omega_{X/B})\longrightarrow\omega_{X/B}
\]
has zero divisor exactly $2S_2$. Consequently
\[
 \omega_X\simeq f^*\mathcal O_B(-p_2)\otimes\mathcal O_X(2S_2)
       \simeq\mathcal O_X(-2S_2),\qquad
 \omega_X^{-2}\simeq f^*\mathcal O_B(p_2).
 \tag{CR1}
\]
For every integer $k$ the actual direct image of the vertical divisor is
\[
 f_*\mathcal O_X(kS_2)
       =\mathcal O_B\bigl(\lfloor k/4\rfloor p_2\bigr)
 \tag{CR2}
\]
as subsheaves of the meromorphic-function sheaf, using pullback to identify
their sections. In particular, for each integer $m\geq0$,
\[
 h^0(X,\omega_X^{-m})=\lfloor m/2\rfloor+1,
 \qquad h^0(X,\omega_X^m)=0\quad(m>0).
 \tag{CR3}
\]
There is a graded algebra isomorphism
\[
 \mathbb C[U,V]\xrightarrow{\sim}
       \bigoplus_{m\geq0}H^0(X,\omega_X^{-m}),
       \qquad\deg U=1,\quad\deg V=2.
 \tag{CR4}
\]
The degree-two complete linear system has no base points and its map is
the original $f$, followed by a specified projective-coordinate map of $B$.
The degree-one system consists of the single divisor $2S_2$.
\end{theorem}



\begin{proof}
\begin{enumerate}[label=\arabic*.,leftmargin=*]
\item Write $\mathfrak K=f_*\omega_{X/B}$. The map is the natural adjunction
evaluation, not an identification inferred from the ranks of its two
terms. On a smooth torus fibre the invariant top form is nowhere zero.
At the cusp, the original form
$(dx_1/x_1)\wedge(dx_2/x_2)\wedge dt_c$ is a nowhere-zero canonical
frame in each of the unimodular toric charts. Its relative version is
the local generator of $\mathfrak K$ used in
\cref{wb:lem:hodge-line-audit}, so evaluation has no zero there.

\item For a finite filling retain $t_j=s_j^{m_j}$,
$(m_1,m_2)=(3,4)$ and $(a_1,a_2)=(2,1)$. On its finite \'etale
total-space cover $q_j:J'_j\to N_j$, the actual direct-image generator
pulls back, by (HL.9.5c)--(HL.9.5f), to
\[
 s_j^{m_j-1-a_j}e^{\varphi_j(s_j)}
       (ds_j\wedge e)\otimes(dt_j)^{-1}.
 \tag{CR5}
\]
Here $(dt_j)^{-1}$ is the pulled-back inverse canonical frame on the
base, not division by $m_js_j^{m_j-1}ds_j$. The coefficient
$e^{\varphi_j}$ is a unit. Thus the zero divisor on the cover is
$(m_j-1-a_j)\{s_j=0\}$. Since $q_j$ is \'etale, faithful pullback
of Cartier divisors makes the downstairs divisor
$(m_j-1-a_j)S_j$. The two coefficients are respectively $0$ and $2$.
These neighbourhoods and the smooth locus cover $X$, proving the stated
global zero divisor with no omitted component.

\item The evaluation therefore identifies
\[
 \omega_{X/B}\simeq f^*\mathfrak K\otimes\mathcal O_X(2S_2).
\]
By \cref{wb:lem:hodge-line-audit}, $\mathfrak K\simeq\mathcal O_B(p_0)$.
In the original affine base coordinate $t$, with
$t(p_1)=0$, $t(p_2)=1$, $t(p_0)=\infty$, retain the explicit choices
\[
 h=\frac1{t-1},\qquad \eta_B=\frac{dt}{(t-1)^2}.
\]
Their divisors are $p_0-p_2$ and $-2p_2$, respectively: at infinity,
with $w=1/t$, they are $w/(1-w)$ and $-dw/(1-w)^2$.
Thus $\eta_B$ specifies $\omega_B\simeq\mathcal O_B(-2p_2)$.
Tensoring with $f^*\omega_B$ gives
\[
 \omega_X\simeq
 \mathcal O_X\bigl(f^*(p_0)-2f^*(p_2)+2S_2\bigr).
\]
Multiplication by $f^*h$, with its displayed divisor, identifies the
base factor $\mathcal O_B(p_0-2p_2)$ with $\mathcal O_B(-p_2)$.
Using the actual multiplicity $f^*(p_2)=4S_2$ proves (CR1).

\item To prove (CR2), a section of $\mathcal O_X(kS_2)$ is represented by
a meromorphic function on $X$ with the indicated vertical order bound.
Away from $p_2$ it is holomorphic. Its restriction to each connected
compact fibre there is constant; $f_*\mathcal O_X=\mathcal O_B$
therefore descends it uniquely to a holomorphic function on that base
open set. Near $p_2$, pull back to the original cover $J'_2$, including
its zero section. The descended function becomes a function $a(s_2)$,
and $s_2^k a(s_2)$ is holomorphic at zero (with the evident vanishing
interpretation when $k<0$). Thus $a$ is meromorphic and has a Laurent
series bounded below. Invariance under the actual quotient, whose base
coordinate is $s_2\mapsto\zeta_2s_2$, says
$a(\zeta_2s_2)=a(s_2)$. Only exponents divisible by four occur, so
$a(s_2)=b(s_2^4)$ for a meromorphic $b$ on the original base disc.
The order inequality is exactly
\[
 4\operatorname{ord}_{p_2}(b)+k\geq0
 \quad\Longleftrightarrow\quad
 \operatorname{ord}_{p_2}(b)\geq-\lfloor k/4\rfloor.
\]
Conversely every such $b$ pulls back to an allowed meromorphic function.
This proves equality in both directions, as well as compatibility on
overlaps. No assumption about the field of all meromorphic functions
on $X$ is used. Applying (CR2) with $k=2m$ and $k=-2m$ and using
$h^0(\mathbb P^1,\mathcal O(d))=\max\{d+1,0\}$ proves (CR3).

\item Here is the exact algebra map in (CR4). Fix the isomorphism
$\omega_X^{-1}\simeq\mathcal O_X(2S_2)$ obtained above. Let $u$ be
the section corresponding to the canonical divisor section of
$\mathcal O_X(2S_2)$, so $\operatorname{div}(u)=2S_2$.
Its square, under the induced degree-two identification in (CR1),
is $f^*\sigma_{p_2}$, where $\sigma_{p_2}$ is the canonical divisor
section of $\mathcal O_B(p_2)$. Choose a second section $\ell$ of this
line bundle with $\ell(p_2)\ne0$, retaining that choice, and let $v$
correspond to $f^*\ell$ under the same degree-two identification.
Define $U\mapsto u$, $V\mapsto v$.

\item For $m=2r$, the monomials $u^{2r-2j}v^j$, $0\leq j\leq r$,
correspond exactly to
$f^*(\sigma_{p_2}^{r-j}\ell^j)$, a basis of
$H^0(B,\mathcal O_B(rp_2))$. For $m=2r+1$, multiplication by $u$
identifies the same basis with the full space: in the meromorphic
description (CR2), both $4rS_2$ and $(4r+2)S_2$ allow precisely
the base pole orders at most $r$. Thus the displayed map is an
isomorphism in every degree, proving both surjectivity and absence of
polynomial relations. Replacing $\ell$ by $a\ell+b\sigma_{p_2}$,
$a\ne0$, changes the second generator by exactly
$v\mapsto av+bu^2$; no scalar is silently discarded.
Replacing the chosen isomorphism $\omega_X^{-1}\to\mathcal O_X(2S_2)$
by $c$ times itself changes the generators by
$(u,v)\mapsto(c^{-1}u,c^{-2}v)$, with $c\in\mathbb C^*$ retained.

\item The two degree-two generators have no common zero because
$\sigma_{p_2},\ell$ do not. Their map is
$x\mapsto[\sigma_{p_2}(f(x)):\ell(f(x))]$, which is exactly
$f$ followed by that chosen coordinate map $B\to\mathbb P^1$.
Explicitly, writing $\ell=A/(t-1)+B$ with $A\ne0$, it is
$t\mapsto[t-1:A+B(t-1)]$, including the values $[0:A]$ at $p_2$
and $[1:B]$ at $p_0$.
Degree one has basis $u$ and its divisor is $2S_2$, as claimed.
\end{enumerate}
\end{proof}

In the already specified additive Picard coordinate \eqref{inv:eq:C-factor12}, with
$L_{12}=f^*\mathcal O_B(1)$, the identities retain the original constants:
\[
 \mathcal O_X(S_1)=L_4,\qquad
 \mathcal O_X(S_2)=L_3,\qquad
 \omega_X=L_{-6},\qquad \omega_X^{-2}=L_{12}.
\]
This last expression uses the previously proved Picard identification;
the divisor, direct-image and ring proofs above do not depend on it.


\subsection{The retained period constant is an effective deformation parameter}
\label{wb:sec:retained-period-deformation}

Keep the functions, lattice marking, base coordinate and local fillings of
\cref{wb:sec:candidate-geometry}. In particular, fix the actual cusp-regular
solution $\beta^{\rm part}$, not a replacement with a cusp pole, and write
\[
 \beta_c=\beta^{\rm part}+c,\qquad
 \mathcal P=\{c\in\mathbb C:\operatorname{Im}c<-M\},
 \quad M=\max_{B\setminus\{p_0\}}D_{\rm part}.
 \tag{PD1}
\]
The maximum and the exact admissible domain were proved in
\cref{wb:prop:period-exact-admissible-constants}. The following calculation
uses the analytic construction, not its integral cohomology or sphere
recognition. Its purpose is to determine whether the freedom in $c$
actually changes the constructed complex manifold.

\begin{theorem}[Exact period deformation]
\label{wb:thm:retained-period-deformation}
\StatementRecord{wb:thm:retained-period-deformation}{proved}{Historical events: GPT-6 Astra (OpenAI), ChatGPT Pro mode, effort Pro, S6\_09\_PAPERS, 17 November 2026: targeted proof check. Also checked in this writing. Pre-continuation writing events: GPT-6 Astra (OpenAI), ChatGPT Pro mode, effort Pro, S6\_10\_PAPERS, 18 November 2026: targeted proof check. Continuation metadata events: GPT-6 Astra (OpenAI), ChatGPT Pro mode, effort Pro, S6\_10\_PAPERS continuation, 19 November 2026: used and rechecked in this run; status-vocabulary reconciliation; metadata reconciliation only; no new mathematical check. \emph{Pre-continuation recorded proof scope.} Read and re-derived the full parameter-retained construction, uniform cusp properness, exact real lattice morphism, plus-graph/connecting-class sign conversion, nonzero torus class, relative canonical-pencil base change over dual numbers, scheme-theoretic constancy of all three critical values, and resulting nonzero Kodaira-Spencer class of the full threefold.}{the hypotheses, notation and definitions stated here; \cref{wb:lem:middle-period-proper}; \cref{wb:lem:toric-chart-control}; \cref{wb:thm:candidate-cusp-geometry}; \cref{wb:thm:canonical-ring-recovery}; \cite[Section~67, p.~15]{McMullen1996Complex}}

For every open disc $Q$ with $\overline Q\subset\mathcal P$, the retained
construction gives a proper holomorphic submersion
\[
 \pi:\mathcal X_Q\longrightarrow Q,
 \qquad F:\mathcal X_Q\longrightarrow B\times Q,
 \qquad \pi=\operatorname{pr}_Q\circ F,
 \tag{PD2}
\]
whose fibre over $c$ is the analytically constructed threefold $X_c$ and
whose restriction there is the original $f_c:X_c\to B$.
The three distinguished values remain the original
$t(p_1)=0$, $t(p_2)=1$, $t(p_0)=\infty$.

Fix a regular base point and one of its original lifts $z$, and fix
$c_*\in Q$. With all functions below evaluated at that $z$, set
\[
 T=\operatorname{Im}\tau>0,\quad m=\operatorname{Im}\mu,\quad
 D=\operatorname{Im}\beta_{c_*}-\frac{6m^2}{T}<0,\quad a=\frac mT,
 \qquad P=\begin{pmatrix}0&0\\-a&1\end{pmatrix}.
 \tag{PD3}
\]
For $\delta=c-c_*$, the exact marked real-linear map from the universal
cover of the fibre at $c_*$ to that at $c$ is
\[
 w=A_\delta\zeta+B_\delta\overline\zeta,\qquad
 A_\delta=I+\frac{\delta}{2iD}P,\qquad
 B_\delta=-\frac{\delta}{2iD}P.
 \tag{PD4}
\]
It maps the original period lattice onto the new period lattice, and
therefore descends to a real-analytic marked diffeomorphism of the two
tori. In the explicit convention that the transported $(0,1)$ tangent
space is spanned by
$\partial_{\overline\zeta_k}+\sum_j\nu^j_k\partial_{\zeta_j}$,
its exact coefficient matrix and its first derivative are
\[
 \nu_\delta=\frac{\delta}{2iD+\delta}P,\qquad
 \dot\nu_z=\frac{1}{2iD}\,
 \partial_{\zeta_2}\otimes
 (d\overline\zeta_2-a\,d\overline\zeta_1).
 \tag{PD5}
\]
The class of $\dot\nu_z$ in
$H^{0,1}(f_{c_*}^{-1}(t(z)),T^{1,0})$ is nonzero.
With the tangent-exact-sequence convention
$\operatorname{KS}(\partial_c)=[\bar\partial h]$ for a smooth
$(1,0)$ lift $h$ of $\partial_c$, the actual fibre class is
$\operatorname{KS}_z(\partial_c)=-[\dot\nu_z]$.
Moreover the first-order deformation of the entire $X_{c_*}$ obtained
from (PD2) is nontrivial: its Kodaira--Spencer class in
$H^1(X_{c_*},T_{X_{c_*}})$ is nonzero. In particular
$h^1(X_{c_*},T_{X_{c_*}})\geq1$.
\end{theorem}



\begin{proof}
\begin{enumerate}[label=\arabic*.,leftmargin=*]
\item \emph{Construction with the parameter retained.}
On an ordinary lifted base open set use the original period map
\[
 \Pi(z,c)=
 \begin{pmatrix}
 6\mu(z)&\tau(z)&1&0\\
 \beta^{\rm part}(z)+c&\mu(z)&0&1
 \end{pmatrix}:
 \Lambda\otimes\mathbb R\longrightarrow\mathbb C^2.
 \tag{PD6}
\]
Its real determinant is $\operatorname{Im}\tau\,D_c$, which never
vanishes on $Q$. Thus $(z,c,u)\mapsto[z,c,\Pi(z,c)u]$ identifies the
quotient, as a real-analytic family, with the original marked real torus
over this base open set times $Q$. The lattice action is holomorphic in
$(z,c,\zeta)$; it is free and properly discontinuous locally uniformly
on compact base-parameter sets, since the inverse real period matrices
are uniformly bounded there. The quotient is therefore a complex
manifold. Projection to $Q$ is a submersion because it is so before
the covering quotient. Its map to the lifted base times $Q$ is proper:
over any compact subset the image of that subset times a fixed closed
lattice parallelotope is the full inverse image.

\item The deck matrices are the original integral $M_g$ and the original
matrices $R_g(z)$, with
$\Pi(gz,c)=R_g(z)\Pi(z,c)M_g$. The $R_g$ do not depend on $c$.
Consequently the ordinary base quotient is jointly holomorphic.
Proper discontinuity there follows from proper discontinuity on the
base, exactly as in \cref{wb:lem:middle-period-proper}, with compact
parameter sets carried along.

\item At a finite filling retain $m_1=3$, $m_2=4$,
$s_j(g_jz)=e^{-2\pi i/m_j}s_j(z)$, and the actual logarithmic lift
\[
 (z,c,\zeta)\longmapsto
 \left(g_jz,c,R_{g_j}(z)\zeta+
                  \Pi(g_jz,c)\frac{v_j}{m_j}\right).
 \tag{PD7}
\]
In the real marking this is $(z,c,u)\mapsto
(g_jz,c,A_ju+v_j/m_j)$, independent of $c$.
Every nonidentity power remains free: the invariant integral covector
$\gamma$ takes the translation of its $k$-th power to $k/3$ or $-k/4$,
which is not an integer for $0<k<m_j$.
The finite quotient is therefore a local biholomorphism and is a
submersion to $Q$. The same compact-parallelotope argument proves
properness over $D_j\times Q$. Its map is still exactly
$t_j=s_j^{m_j}$. On the punctured overlap the translating section is
\[
 \frac{\log s_j(z)}{2\pi i}\Pi(z,c)v_j.
 \tag{PD8}
\]
A change of logarithm adds an actual period. It is therefore a
branch-independent holomorphic torus section, jointly in $z,c$.
The calculation (2.9e) applies with $c$ retained and proves the
punctured-overlap conjugacy, not merely equality of the monodromy orders.

\item For the cusp retain the same infinite unimodular fan and the full matrix
\[
 C(t_c,c)=
 \begin{pmatrix}6\mu(t_c)&h(t_c)\\
 b^{\rm part}(t_c)-h(t_c)+c&\mu(t_c)
 \end{pmatrix},\qquad
 B_0=\begin{pmatrix}0&1\\-1&0\end{pmatrix}.
 \tag{PD9}
\]
Choose a closed cusp disc of radius $r_1$ on which these functions are
holomorphic. Put
\[
 K_Q=\max_{|t_c|\le r_1,\ c\in\overline Q}
          \|-2\pi\operatorname{Im}C(t_c,c)\|,
 \qquad 0<\epsilon<\min(r_1,1),\quad
          |\log\epsilon|>2K_Q.
 \tag{PD10}
\]
The action on $Y_\epsilon\times Q$ is the original one,
\[
 (x,t_c,c)\longmapsto
 \left(e^{2\pi iC(t_c,c)\lambda}
                  t_c^{B_0\lambda}x,t_c,c\right),
 \quad\lambda\in\mathbb Z^2.
 \tag{PD11}
\]
It is holomorphic across the boundary by the same toric automorphism
and nowhere-zero torus multiplier. The estimates in (2.8) now hold
uniformly on $Q$:
\[
 B_{t_c,c}=B_0+
       \frac{-2\pi\operatorname{Im}C(t_c,c)}{\log|t_c|},\quad
 \|B_{t_c,c}-B_0\|<\tfrac12,\quad
 \|B_{t_c,c}\lambda\|>\tfrac12\|\lambda\|\quad(\lambda\ne0).
 \tag{PD12}
\]
They exclude stabilizers on the dense torus and bound the finite set
of translations that can join each pair of bounded charts, uniformly
in $c$. At $t_c=0$ a nonzero translation moves the bounded cell, so
there is no stabilizer there either. A compact set is contained in finitely
many of the open charts $\mathcal U_\sigma(2,\epsilon)\times Q$ from
\cref{wb:lem:toric-chart-control}.
The chart-pair estimate proves proper discontinuity on the whole
product, including sequences tending to its boundary.

\item For clarity properness is also uniform, not only pointwise in $c$.
For $\epsilon_1<\epsilon$, let $K_0$ be the finite union of closed
unit polydiscs associated to triangles meeting the closed radius
$\sqrt2$ ball, intersected with $|t_c|\le\epsilon_1$, exactly as in
\cref{wb:thm:candidate-cusp-geometry}. For every $c$ and $t_c\ne0$,
choose $\lambda$ so that
$\|B_{t_c,c}^{-1}y+\lambda\|_\infty\le1/2$.
The translated logarithmic position has norm at most
$3/(2\sqrt2)<\sqrt2$, and hence lies in a chart defining $K_0$.
At $t_c=0$, the six closed bidiscs covering $E_0$ give the same
representatives after the exact translation by $B_0v$, with multiplier
$\exp(2\pi iC(0,c)B_0v)$ retained. Therefore, for compact $K\subset Q$,
the image of the compact $K_0\times K$ is the full inverse image of
$\{|t_c|\le\epsilon_1\}\times K$. Every compact set in the cusp
base times $Q$ is closed in a set of this form. This proves properness.
The quotient is a submersion to $Q$ because its covering charts are
open subsets of $Y_\epsilon\times Q$.

\item Exponentiating the original logarithmic torus coordinates gives the
jointly holomorphic cusp overlap, with period identity
$Z=sB_0+C(t_c,c)$. Glue it and (PD8) to the middle family.
The special base discs are disjoint. Each resulting inverse image of
a member of this base cover is its entire Hausdorff local piece.
Two points with different $(t,c)$ values are separated on the base;
two with the same value lie in one Hausdorff open piece and can be
separated there. This proves Hausdorffness. Local complex charts give
a complex fourfold and the submersion $\pi$.
Properness of $F$ follows by covering any compact set of $B\times Q$
by finitely many relatively compact sets inside these base pieces and
using the just proved local properness. Since $B$ is compact, projection
$B\times Q\to Q$ is proper, proving (PD2). The fibres and overlaps
are precisely those in the original analytic construction.

\item \emph{The exact variation on a regular fibre.}
Write $u=(u_1,u_2,u_3,u_4)^t$ in the original ordered lattice basis and
$\zeta=\Pi(z,c_*)u$. Taking imaginary parts gives the two exact equations
\[
 \operatorname{Im}\zeta_1=6m u_1+T u_2,\qquad
 \operatorname{Im}\zeta_2=(D+6m^2/T)u_1+m u_2.
 \tag{PD13}
\]
Subtracting $a$ times the first equation from the second yields
\[
 u_1=\frac{\operatorname{Im}\zeta_2-a\operatorname{Im}\zeta_1}{D}.
 \tag{PD14}
\]
In the changed periods the same $u$ maps to
$w=\Pi(z,c)u=\zeta+\delta e_2u_1$. Substitution of
$\operatorname{Im}\zeta=(\zeta-\overline\zeta)/(2i)$ proves (PD4).
This also proves its asserted lattice map and real invertibility:
it is exactly $\Pi(z,c)\Pi(z,c_*)^{-1}$ as a real-linear map.
Its real determinant is $(D+\operatorname{Im}\delta)/D>0$.
Neither the nonholomorphic part nor the original marking is dropped.

\item Here $P^2=P$, and
\[
 \det A_\delta=1+\frac\delta{2iD},\qquad
 A_\delta^{-1}=I-\frac{\delta}{2iD+\delta}P.
 \tag{PD15}
\]
The denominator can vanish only for $\delta=-2iD$. At that parameter
$D+\operatorname{Im}\delta=-D>0$, which is outside the admissible
domain. Thus this inverse is defined at every $c\in Q$.
Annihilating $w$ by the declared new $(0,1)$ fields gives
$A_\delta\nu_\delta+B_\delta=0$, which proves (PD5), including
its sign. All its coefficients are constant in the fibre coordinates.
Thus its first derivative is a globally defined $\bar\partial$-closed
tangent-valued $(0,1)$ form on the torus.

\item It cannot be $\bar\partial V$ for a globally defined smooth vector
field $V$. Indeed the invariant holomorphic frame
$\partial_{\zeta_1},\partial_{\zeta_2}$ writes $V$ with periodic
smooth coefficients. The integral of each of their derivatives
$\partial_{\overline\zeta_k}V^j$ against the translation-invariant
probability volume is zero: real constant vector fields preserve this
volume, and integration of their derivatives on the compact torus is
zero. But the coefficient of
$\partial_{\zeta_2}\otimes d\overline\zeta_2$ in $\dot\nu_z$ is
the nonzero constant $1/(2iD)$. Its integral is that same constant.
This contradiction proves nonzero Dolbeault class without a dimension
count or an appeal to numerical period variation.

\item The conversion to the tangent-exact-sequence convention retains a
minus sign. In the moving coordinates the real-marking lift has
$(1,0)$ part $h=\partial_c+u_1\partial_{w_2}$ at $c_*$.
Since (PD14) gives
\[
 \bar\partial u_1=-\frac{d\overline\zeta_2-a\,d\overline\zeta_1}{2iD},
\]
the connecting class $[\bar\partial h]$ is exactly $-[\dot\nu_z]$.
Both conventions therefore give a nonzero class; they are not identified
without this sign conversion.

\item \emph{The intrinsic pencil in the parameter family.}
Let $\mathcal S_2\subset\mathcal X_Q$ be the reduced finite divisor
over $p_2\times Q$, so $F^*(p_2\times Q)=4\mathcal S_2$.
The proof of \cref{wb:thm:canonical-ring-recovery} works for relative
forms over $Q$, as can be checked on each original chart.
On the middle family the frame is $d\zeta_1\wedge d\zeta_2$ relative
to $B\times Q$. At the cusp the relative canonical frame over $Q$ is
$(dx_1/x_1)\wedge(dx_2/x_2)\wedge dt_c$.
Extra $dc$ terms in transition differentials vanish in relative forms.
The finite determinant characters $\chi_j$, their correcting units
$e^{\varphi_j(s_j)}$, and the modular function $\Theta$ used in
(HL.9.4)--(HL.9.6) are independent of $c$, since $R_g,\tau,\mu$ are.
Thus their coefficient identities hold jointly over $Q$.
The evaluation exponents remain $m_j-1-a_j=0,2$.
It follows by the actual same transition functions and divisors that
\[
 \omega_{\mathcal X_Q/Q}\simeq\mathcal O_{\mathcal X_Q}(-2\mathcal S_2),
 \qquad
 \omega_{\mathcal X_Q/Q}^{-2}\simeq
          F^*\operatorname{pr}_B^*\mathcal O_B(p_2).
 \tag{PD16}
\]
Here the base identifications retain $h=1/(t-1)$ and
$\eta_B=dt/(t-1)^2$ from (CR1); these functions also do not depend on $c$.
In particular (PD16) is a relative line-bundle isomorphism, not a
collection of separately chosen isomorphisms of its fibres.

\item We need its complete degree-two system after an infinitesimal base
change, and prove this explicitly. Put
$A=\mathbb C[\epsilon]/(\epsilon^2)$, base-change by
$c=c_*+\epsilon$, write $X_A$ for the resulting deformation and
$L_A=\omega_{X_A/A}^{-2}$. For $L_*=\omega_{X_{c_*}}^{-2}$,
flatness of the deformation and local freeness of $L_A$ give
\[
 0\longrightarrow L_*\xrightarrow{\ \epsilon\ }L_A
        \longrightarrow L_*\longrightarrow0.
 \tag{PD17}
\]
The two sections of $\mathcal O_B(p_2)$ used in (CR4), namely
$\sigma_{p_2}$ and $\ell=A_0/(t-1)+B_0'$ with $A_0\ne0$,
pull back by (PD16) to sections $s_0,s_1$ of $L_A$ whose reductions
are a basis of $H^0(X_{c_*},L_*)$. This basis assertion is exactly
the proved degree-two case of (CR2)--(CR4).
Any section of $L_A$ reduces to a complex linear combination of these
two reductions; subtract that lifted combination. Exactness of (PD17)
then writes the remainder as $\epsilon$ times a section of $L_*$,
again a combination of that basis. Conversely, if an $A$-linear
combination of $s_0,s_1$ is zero, reduce modulo $\epsilon$ and then
use the injective first arrow of (PD17) to see that all four complex
coefficients vanish. Therefore
\[
 H^0(X_A,L_A)=A s_0\oplus A s_1.
 \tag{PD18}
\]
The system is basepoint-free and its morphism is $f_A$ followed by
the fixed coordinate map $t\mapsto[t-1:A_0+B_0'(t-1)]$.
This proves the required base-change statement without presuming that
pointwise dimensions alone imply it.

\item \emph{Nontriviality for the complete threefold.}
Suppose its first-order deformation were trivial, with an isomorphism
$X_A\simeq X_{c_*}\times\operatorname{Spec}A$ reducing to the identity.
Intrinsic relative canonical bundles and their complete spaces of
sections are preserved by that isomorphism. By (PD18), it would
therefore identify $f_A$ with $a_A\circ f_{c_*}$ for a projective
linear change $a_A\in\operatorname{PGL}_2(A)$ reducing to the identity.

\item This change fixes all three distinguished points even to first order.
Here is a scheme-sensitive check. In the finite covering charts
$f_A$ is $t_j=s_j^{m_j}$, so its relative critical ideal is
$(s_j^{m_j-1})$, and its restriction to that critical scheme has
$t_j=0$. The map from the parameter ring $A$ to this critical
scheme is injective; it contains a constant-parameter section given
by $s_j=0$ and the zero section of the covering torus.
At the cusp, in a chart containing a triple point,
$t_c=z_0z_1z_2$ and the critical ideal is
$(z_1z_2,z_0z_2,z_0z_1)$. Again $t_c=0$ in its ring and $A$ injects,
using $z_0=z_1=z_2=0$. At a double chart the same assertion follows
from $t_c=z_0z_1$; a smooth chart has no critical point.
These critical ideals are intrinsic to the differential of the map
over $A$, and the charts are local biholomorphisms onto $X_A$.
Away from the three distinguished base values the map is a
submersion, so it has no critical point.
Consequently the scheme-theoretic critical-value sections are exactly
the three constant sections $p_1,p_2,p_0$. An automorphism reducing
to the identity cannot permute these three distinct closed points.
Its value at each is therefore the same point over $A$, not a
nonzero infinitesimal displacement: such a displacement would be a
nonzero element of $\epsilon\mathbb C$ in an injective copy of $A$
on the just displayed critical chart.

\item For completeness write $a_A(t)=t+\epsilon(b_0+b_1t+b_2t^2)$.
Fixing $0$ gives $b_0=0$, fixing $1$ gives $b_0+b_1+b_2=0$,
and fixing infinity, in $u=1/t$, gives $b_2=0$.
Hence $a_A$ is the identity. The hypothesized trivialization thus
preserves $f$, and restricting it to any regular base value would
trivialize the torus deformation (PD4). Its nonzero class (PD5)
excludes this.

\item Finally, the zero Kodaira--Spencer class is precisely the condition for
such a first-order trivialization. One may see the needed direction
directly: choose local product charts for the holomorphic submersion
$\pi$ and differentiate their transition functions at $c_*$. The
resulting tangent-valued overlap functions form a Cech $1$-cocycle.
If its class vanishes, after refinement it is the coboundary of local
holomorphic vector fields. Changing each chart by that vector field
times $\epsilon$ cancels the first-order overlap terms, producing
the identity-reducing trivialization. This is the usual
Kodaira--Spencer connecting class (for the standard deformation-space
interpretation see \cite[Section~67, p.~15]{McMullen1996Complex});
the explicitly recorded conversion
from the plus-graph convention in (PD5) changes its sign, not its
vanishing. The contradiction proves the final assertion.
\end{enumerate}
\end{proof}

The theorem establishes a genuinely varying analytic family with its
original marked fibration and its intrinsic degree-two canonical pencil.
It does not assert that every two distinct parameters give nonisomorphic
unmarked manifolds, or that this family exhausts their deformation space.
The global integral topology and the CDP argument are not inputs to this
deformation calculation.


The global identifications within this retained family are computed in
\cref{wb:pshift:thm:classification}: they are precisely $c'-c\in2\mathbb Z$,
including the finite-filling obstruction and every cusp overlap.


\subsection{Exact global identifications of the retained period parameter}
\label{wb:pshift:sec:classification}

Keep the original base $B$, coordinate $t$, values
$t(p_1)=0,t(p_2)=1,t(p_0)=\infty$, and the original functions
$\tau,\mu,\beta^{\mathrm{part}}$. Write
\[
 \beta_c=\beta^{\mathrm{part}}+c,\qquad
 \mathcal P=\{c\in\mathbb C:\operatorname{Im}c<-M\},\qquad
 \Pi_c(z)=\begin{pmatrix}6\mu(z)&\tau(z)&1&0\\
                      \beta_c(z)&\mu(z)&0&1\end{pmatrix}.
 \tag{PS1}\label{wb:pshift:eq:periods}
\]
All assertions below concern the analytic pieces and their actual gluing.
No integral cohomology, sphere recognition, or CDP assertion is an input.
In particular, the global parameter classification is not inferred from
an equality between two fibre lattices.

Put $e=(0,1)^{\mathsf t}\in\mathbb C^2$ and retain the ordered lattice
\[
 \Lambda=\mathbb Z\langle\widehat\gamma,\widehat u,
                        \widehat w,\widehat\delta\rangle,
 \qquad E=\widehat\delta\otimes\gamma.
 \tag{PS2}\label{wb:pshift:eq:rank-one-map}
\]
Thus $E$ is the matrix with its only nonzero entry, equal to one, in
row four and column one; $E^2=0$ and $\Pi_cE=e\gamma$.
The finite data, with every sign retained, are
\[
\begin{gathered}
 A_1=\begin{pmatrix}1&0&0&0\\6&0&1&0\\-6&-1&-1&0\\-2&1&0&1\end{pmatrix},
 \qquad
 A_2=\begin{pmatrix}1&0&0&0\\0&0&-1&0\\-6&1&0&0\\3&0&1&1\end{pmatrix},\\
 v_1=(1,2,-4,0)^{\mathsf t},\quad
 v_2=(-1,-3,3,0)^{\mathsf t},\quad
 (m_1,m_2)=(3,4),\quad
 (\epsilon_1,\epsilon_2)=(1,-1),\quad \epsilon_j=\gamma(v_j).
\end{gathered}
\tag{PS3}\label{wb:pshift:eq:finite-data}
\]
The covariance is $R_g\Pi_c=\Pi_c(gz)A_g$, with
\[
 R_{g_1}=\begin{pmatrix}-1/\tau&0\\(1-\mu)/\tau&1\end{pmatrix},\qquad
 R_{g_2}=\begin{pmatrix}1/\tau&0\\-\mu/\tau&1\end{pmatrix},\qquad
 R_{g_0}=I_2.
 \tag{PS4}\label{wb:pshift:eq:covariance}
\]
In particular $R_ge=e$ for every $g$.

We give the full period-shift proof of Theorem~\ref{fam:st:A.5.1}.


\subsubsection{The middle family and its complete rational centralizer}

For an integer $n$, put $S_n=I+nE$. Since $S_n^{-1}=S_{-n}$,
\[
 \Pi_{c+n}=\Pi_c S_n,\qquad
 \Pi_{c+n}\Lambda=\Pi_c\Lambda.
 \tag{PS6}\label{wb:pshift:eq:integer-shear}
\]
Both original matrices fix $\widehat\delta$ and preserve $\gamma$.
Consequently $A_jE=EA_j=E$, and the same commutation holds for every
monodromy word. The identity of $\mathbb C^2$ therefore identifies
the two full middle torus families with their actual $R_g$-action.
Its source-to-target lattice map is $S_{-n}$, since
$\Pi_{c+n}S_{-n}=\Pi_c$.

\begin{lemma}[All middle-family linear parts]
\label{wb:pshift:lem:centralizer}
\StatementRecord{wb:pshift:lem:centralizer}{proved}{Historical events: GPT-6 Astra (OpenAI), ChatGPT Pro mode, effort Pro, S6\_09\_PAPERS, 17 November 2026: targeted proof check. Also checked in this writing. Pre-continuation writing events: GPT-6 Astra (OpenAI), ChatGPT Pro mode, effort Pro, S6\_10\_PAPERS, 18 November 2026: targeted proof check. Continuation metadata events: GPT-6 Astra (OpenAI), ChatGPT Pro mode, effort Pro, S6\_10\_PAPERS continuation, 19 November 2026: used and rechecked in this run; status-vocabulary reconciliation; metadata reconciliation only; no new mathematical check. \emph{Pre-continuation recorded proof scope.} Re-derived the entire simultaneous rational centralizer using both fixed vector and fixed covector planes and the two quotient matrices; solved the extension equations, imposed integrality and determinant, and compared all period columns to obtain the integral shift condition.}{the hypotheses, notation and definitions stated here; the explicit constructions and calculations in the proof}

The simultaneous rational centralizer of $A_1,A_2$ is exactly
\[
 \{aI+bE:a,b\in\mathbb Q\}.
 \tag{PS7}\label{wb:pshift:eq:centralizer}
\]
A biholomorphism between the middle families for $c,c'$ over the
identity of $B^\circ$ forces $c'-c\in\mathbb Z$. Its complex-linear
part on every fibre is either $I_2$ or $-I_2$.
\end{lemma}



\begin{proof}
\begin{enumerate}[label=\arabic*.,leftmargin=*]
\item Solving the two fixed-vector systems from (PS3) gives
\[
 (\Lambda\otimes\mathbb Q)^{A_1}
   =\mathbb Q(1,2,-4,0)^{\mathsf t}+\mathbb Q\widehat\delta,
 \quad
 (\Lambda\otimes\mathbb Q)^{A_2}
   =\mathbb Q(1,3,-3,0)^{\mathsf t}+\mathbb Q\widehat\delta.
\]
Their intersection is $\mathbb Q\widehat\delta$. The fixed covectors
are respectively
\[
 \mathbb Q\gamma+\mathbb Q(2u+w+3\delta),\qquad
 \mathbb Q\gamma+\mathbb Q(u+w+2\delta).
\]
Their intersection is $\mathbb Q\gamma$: comparing the $u$ and $w$
coefficients of a common multiple gives $2x=y$ and $x=y$, hence
$x=y=0$. Every commuting endomorphism therefore preserves
$\mathbb Q\widehat\delta$ and $\ker\gamma$.

\item On $\ker\gamma/\mathbb Q\widehat\delta$, in the classes of
$(\widehat u,\widehat w)$, the two matrices are
\[
 Q_1=\begin{pmatrix}0&1\\-1&-1\end{pmatrix},\qquad
 Q_2=\begin{pmatrix}0&-1\\1&0\end{pmatrix}.
\]
A matrix commuting with $Q_2$ has the form
$\left(\begin{smallmatrix}b&d\\-d&b\end{smallmatrix}\right)$.
Commutation with $Q_1$ gives $d=0$. Thus a commuting endomorphism $H$
has columns
\[
 H\widehat\gamma=(a,r,s,k)^{\mathsf t},\quad
 H\widehat u=b\widehat u+u_0\widehat\delta,\quad
 H\widehat w=b\widehat w+w_0\widehat\delta,\quad
 H\widehat\delta=d_0\widehat\delta.
\]
Use the actual columns $A_1\widehat u=-\widehat w+\widehat\delta$,
$A_1\widehat w=\widehat u-\widehat w$, and
$A_2\widehat u=\widehat w$. Commutation on these columns gives
\[
 d_0-w_0=b+u_0,\qquad u_0=2w_0,\qquad w_0=u_0.
\]
Hence $u_0=w_0=0$ and $d_0=b$. Commutation on
$\widehat\gamma$ with $A_1$ gives
\[
 r=2(a-b),\qquad s=-4(a-b).
\]
Commutation on that column with $A_2$ also gives $r=-s$.
Therefore $a=b$, $r=s=0$, and $H=aI+kE$. Conversely these matrices
commute, by the fixed-vector and fixed-covector identities preceding
(PS6). This proves (PS7).

\item For the geometric assertion, a holomorphic map between two compact
complex tori lifts to their universal covers. The derivative of the
lift is periodic, so its entries descend to holomorphic functions on
the compact source torus and are constant. The lift is consequently
affine, with a complex-linear part $L$. In a holomorphic family these
linear parts vary holomorphically. Their lattice maps are integral
and locally constant, hence give an isomorphism of the two marked
local systems. Because the map is over the identity of the base,
that integral map $H$ commutes with both $A_j$. Thus
$H=aI+bE$, with $a,b\in\mathbb Z$. Its determinant is $a^4$;
invertibility over $\mathbb Z$ gives $a=1$ or $a=-1$.

\item The complete period equation is
\[
                   L(z)\Pi_c(z)=\Pi_{c'}(z)(aI+bE).
 \tag{PS8}\label{wb:pshift:eq:period-morphism}
\]
Its third and fourth columns give $L(z)=aI_2$. Its first column
then gives $a(c'-c)+b=0$. Therefore $c'-c=-b/a\in\mathbb Z$,
and the asserted two linear parts exhaust the possibilities.
\end{enumerate}
\end{proof}

\subsubsection{The exact finite obstruction, including normalization}

The two cyclic norm matrices are, by multiplication of (PS3),
\[
 N_1=I+A_1+A_1^2=
 \begin{pmatrix}3&0&0&0\\6&0&0&0\\-12&0&0&0\\0&2&1&3\end{pmatrix},
\]
\[
 N_2=I+A_2+A_2^2+A_2^3=
 \begin{pmatrix}4&0&0&0\\12&0&0&0\\-12&0&0&0\\0&2&2&4\end{pmatrix}.
 \tag{PS9}\label{wb:pshift:eq:norms}
\]
For example these formulas give
$N_1\widehat w=\widehat\delta$,
$N_2\widehat u=N_2\widehat w=2\widehat\delta$,
$N_1\widehat\delta=3\widehat\delta$, and
$N_2\widehat\delta=4\widehat\delta$.

We first justify why a local isomorphism must satisfy the exact affine
equations on the original finite cover. The actual ramified base change
$t_j=u^{m_j}$ has local equation
\[
 s_j^{m_j}-u^{m_j}
   =\prod_{a=0}^{m_j-1}(s_j-e^{2\pi ia/m_j}u)=0.
\]
Its normalization separates all these smooth branches, each with
coordinates $(z_1,z_2,u)$. The original map
$y\mapsto(q_j(y),s_j(y))$ identifies the original torus family $J'_j$
with this normalization, as proved by the explicit branches in
\cref{wb:ffcert:thm:normalization}. A biholomorphism over the identity of
$D_j$ base-changes over the identity of the $u$-disc and lifts uniquely
to the normalization: on local rings it carries integral elements in
the total fraction ring to integral elements, and its inverse does
the same. This lift commutes with the original deck action because
both composites lift the same map $(x,u)\mapsto(F(x),\zeta_j u)$.
In particular one may not replace the generator by an arbitrary
coprime power.

Let $c'=c+n$, where $n$ is an integer, and suppose the linear part of
the lifted map is $aI_2$, $a\in\{1,-1\}$. The lifted map on the torus
family is $[z,\zeta]\mapsto[z,a\zeta+B_j(z)]$. Its translating
section has a holomorphic lift $B_j:\Delta_j\to\mathbb C^2$:
$\Delta_j\times\mathbb C^2\to J'_j$ is a covering map and the disc is
simply connected. Local inverses of this covering show that the lift
is holomorphic. The deck-equivariance identity holds modulo an
actual lattice vector, so, after choosing this lift, it is
\[
 B_j(g_jz)-R_{g_j}(z)B_j(z)
  =\Pi_c(g_jz)\left(
       \frac{(1-a)v_j+n\epsilon_j\widehat\delta}{m_j}
          +\ell_j\right),\qquad \ell_j\in\Lambda.
 \tag{PS10}\label{wb:pshift:eq:finite-equation}
\]
The integer column $\ell_j$ is constant: it is locally constant by
discreteness of the marked lattice and the disc is connected.
Iterating $m_j$ times, with $A_jv_j=v_j$ and
$A_j\widehat\delta=\widehat\delta$, yields the necessary equality
\[
 N_j\ell_j=-(1-a)v_j-n\epsilon_j\widehat\delta.
 \tag{PS11}\label{wb:pshift:eq:finite-norm-obstruction}
\]
Here the left side of the iterated equation is zero because the
base point and the original vector coordinate return after $m_j$
steps; real invertibility of $\Pi_c$ makes the lattice equality exact.

If $a=-1$, applying $\gamma$ to (PS11) gives
\[
                   m_j\gamma(\ell_j)=-2\epsilon_j.
\]
This is impossible already at $j=1$, since $3$ does not divide $2$.
Thus a global isomorphism cannot have negative linear part.

For $a=1$, the first row of (PS11) gives $\gamma(\ell_j)=0$.
Equation (PS9) then proves that the order-three filling imposes no
restriction on $n$, whereas the order-four filling imposes precisely
$n\in2\mathbb Z$. More intrinsically, the finite obstruction is the
class of $-n\epsilon_j\widehat\delta$ in
\[
 \Lambda^{A_j}/N_j\Lambda.
 \tag{PS12}\label{wb:pshift:eq:finite-class}
\]
In the original fixed bases this quotient is
$\mathbb Z/3$ at $j=1$, with $\widehat\delta$ zero, and
$\mathbb Z/4\oplus\mathbb Z/2$ at $j=2$, with
$\widehat\delta$ generating the second summand. Indeed (PS9) gives
\[
 N_1\Lambda=3\mathbb Z v_1\oplus\mathbb Z\widehat\delta,
 \quad
 N_2\Lambda=4\mathbb Z\varepsilon'\oplus
                      2\mathbb Z\widehat\delta,
 \qquad \varepsilon'=(1,3,-3,0)^{\mathsf t}.
\]
This proves the exact parity obstruction rather than merely failure
of the untranslated identity.

\subsubsection{Local maps for every even shift}

Let $n\in2\mathbb Z$. Define the following original rational columns,
integral cocycles, and holomorphic translation vectors:
\[
\begin{aligned}
 r_1&=-\frac n3(\widehat u+\widehat w),&
 \ell_1&=-n\widehat w,&
 b_1(z)&=\Pi_c(z)r_1=-\frac n3(\tau(z)+1,\mu(z))^{\mathsf t},\\
 r_2&= \frac n4(\widehat u+\widehat w),&
 \ell_2&=\frac n2\widehat u,&
 b_2(z)&=\Pi_c(z)r_2=\frac n4(\tau(z)+1,\mu(z))^{\mathsf t}.
\end{aligned}
\tag{PS13}\label{wb:pshift:eq:local-corrections}
\]
Direct multiplication of the complete matrices in (PS3) gives
\[
 (I-A_j)r_j=\frac{n\epsilon_j}{m_j}\widehat\delta+\ell_j.
 \tag{PS14}\label{wb:pshift:eq:local-matrix-identity}
\]
For $j=1$ the four entries are $(0,0,-n,n/3)$, and for $j=2$
they are $(0,n/2,0,-n/4)$. Pairing with the full $\Pi_c(g_jz)$
shows that $b_j$ satisfies (PS10) for $a=1$.
Translation by $b_j$ is therefore already an isomorphism of the
two finite local fillings. The remaining issue is its exact overlap
with the fixed middle identification.

Keep the original linearizing coordinate, radius, and phase
$s_j(g_jz)=e^{-2\pi i/m_j}s_j(z)$. On a logarithm branch put
\[
 w_j(z)=\frac{\log s_j(z)}{2\pi i},\qquad
 a_j(z)=b_j(z)+n\epsilon_j w_j(z)e.
 \tag{PS15}\label{wb:pshift:eq:finite-middle-sections}
\]
Changing the logarithm adds $n\epsilon_j\widehat\delta$ under the
period map, so $a_j$ is branch independent as a torus section.
For the specified negative generator, $w_j(g_jz)=w_j(z)-1/m_j$.
Equations (PS4) and (PS14) consequently give the exact equality
\[
 a_j(g_jz)=R_{g_j}(z)a_j(z)+\Pi_c(g_jz)\ell_j.
 \tag{PS16}\label{wb:pshift:eq:local-affine-cocycle}
\]
This is why the logarithmic term in (PS15) cannot be omitted.
In fact the original finite overlap maps are translations by
$\sigma_{j,c}=w_j\Pi_c v_j$, and
\[
                  \sigma_{j,c+n}-\sigma_{j,c}
                     =n\epsilon_jw_je.
 \tag{PS17}\label{wb:pshift:eq:actual-overlap-difference}
\]
Thus $a_j$ is exactly the middle translation induced by the local
translation $b_j$, with the original gluing maps retained.

At the cusp retain $s=\tau-h$, $e^{2\pi is}=t_c$, $s(g_0z)=s(z)-1$,
and the complete period matrix
\[
 Z_c=sB_0+C_c(t_c),\quad
 B_0=\begin{pmatrix}0&1\\-1&0\end{pmatrix},\quad
 C_c=\begin{pmatrix}6\mu&h\\ b^{\mathrm{part}}-h+c&\mu\end{pmatrix}.
 \tag{PS18}\label{wb:pshift:eq:cusp-periods}
\]
For $\lambda=(\lambda_1,\lambda_2)^{\mathsf t}\in\mathbb Z^2$,
$C_{c+n}\lambda-C_c\lambda=(0,n\lambda_1)^{\mathsf t}$. Therefore
\[
 \exp(2\pi iC_{c+n}\lambda)t_c^{B_0\lambda}
       =\exp(2\pi iC_c\lambda)t_c^{B_0\lambda}.
 \tag{PS19}\label{wb:pshift:eq:cusp-actions-identical}
\]
The entire cusp actions, including their nonconstant holomorphic units,
are identical for an integer shift.

The affine cocycle convention in (PS16) has multiplication law
\[
                    \ell_{gh}=\ell_g+A_g\ell_h.
 \tag{PS20}\label{wb:pshift:eq:cocycle-law}
\]
Since $g_0=(g_1g_2)^{-1}$, it gives
\[
\begin{aligned}
 \ell_1+A_1\ell_2
   &=-\frac{3n}{2}\widehat w+\frac n2\widehat\delta,\\
 \ell_0=-A_{g_0}(\ell_1+A_1\ell_2)
   &=\frac{3n}{2}\widehat w-\frac n2\widehat\delta.
\end{aligned}
\tag{PS21}\label{wb:pshift:eq:cusp-cocycle}
\]
Here $A_{g_0}$ fixes both toric lattice vectors, as is also immediate
from $A_{g_0}=\left(\begin{smallmatrix}I&0\\B_0&I\end{smallmatrix}\right)$.
Put
\[
 q_n=(-3n/2,n/2)^{\mathsf t}\in\mathbb Z^2,
                  \qquad a_0(z)=s(z)q_n.
 \tag{PS22}\label{wb:pshift:eq:cusp-translation}
\]
Then $a_0(g_0z)-a_0(z)=-q_n=\Pi_c(g_0z)\ell_0$, exactly as
required by (PS16). Exponentiating gives the map
$(x,t_c)\mapsto(t_c^{q_n}x,t_c)$ on the dense cusp torus.
It extends holomorphically and invertibly to the original infinite fan:
the integral map
\[
 (y,y_3)\longmapsto(y+y_3q_n,y_3)
\]
translates each height-one triangle by the integer vector $q_n$ and
therefore permutes the actual cones, with inverse given by $-q_n$.
It commutes with the cusp quotient action in (PS19). This proves
extension across every cusp stratum, not only across its dense torus.

\subsubsection{The global additive correction and all overlap identities}

We give the global correction explicitly as an additive gluing
construction and prove its solvability. Let $\widetilde B^\circ$ be
the universal cover of the three-punctured base. Its fundamental group
is free on the original two negative peripheral generators. Pull back
the original $\Pi_c,R$ to this cover. Assign $\ell_1,\ell_2$ from
(PS13) to its free generators and extend by (PS20). Freeness makes
this an actual integral cocycle; its third peripheral value is
precisely (PS21). The maps
\[
 (z,y)\longmapsto
       (gz,R_g(z)y+\Pi_c(gz)\ell_g)
 \tag{PS23}\label{wb:pshift:eq:affine-bundle-action}
\]
form an affine action. To check the action law, covariance gives
$R_g(hz)\Pi_c(hz)=\Pi_c(ghz)A_g$; substituting this equality
in the composition gives exactly (PS20).
The quotient is an affine holomorphic bundle over $B^\circ$, since
the deck action on the base is free and properly discontinuous and
one may use evenly covered base opens as affine charts.

On a finite punctured disc, (PS16) says that $a_j$ is a section of
this affine bundle. On the cusp punctured disc, (PS21)--(PS22) say
the same for $a_0$. Define a sheaf $\mathscr H_n$ on $B$ as follows.
Away from the special points it is the sheaf of sections of (PS23).
Near $p_j$, a section is required to have the form $a_j+\eta_j$,
where $\eta_j:\Delta_j\to\mathbb C^2$ is holomorphic and obeys
$\eta_j(g_jz)=R_{g_j}(z)\eta_j(z)$. Near the cusp it is required
to have the form $a_0+\eta_0$, with $\eta_0(t_c)$ holomorphic on
the full cusp disc. These are compatible definitions on the punctured
overlaps by the equal affine cocycles just proved. They also prove
local nonemptiness, including at all three special points.

The differences of these local sections form a vector sheaf
$\mathscr V$. Its ordinary sections have covariance $\eta(gz)=R_g\eta(z)$;
at the two finite points and the cusp its extensions are precisely
the holomorphic extensions specified above. Projection to the first
coordinate gives an exact sequence on the original compact base
\[
 0\longrightarrow\mathcal O_B e\longrightarrow\mathscr V
        \longrightarrow\mathcal L\longrightarrow0,
                  \qquad \mathcal L\simeq\mathcal O_B(-p_0).
 \tag{PS24}\label{wb:pshift:eq:vector-sheaf}
\]
Here $\mathcal L$ is exactly the homogeneous period sheaf from
\cref{wb:prop:period-input-audit}, not a newly chosen extension. Indeed
the first diagonal entries in (PS4) are its factors $-1/\tau$ and
$1/\tau$, and at the cusp both extensions require ordinary
holomorphicity. The original function
\[
 \Theta=\frac{(E_4\circ\tau)^2(E_6\circ\tau)^{1/2}}
                 {\Delta_{\mathrm{mod}}\circ\tau}
\]
has those exact factors, orders two and one in $s_1,s_2$, and a
simple pole in $t_c$. Dividing by $\Theta$ identifies the ordinary
and finite stalks with $\mathcal O_B$, and the cusp stalk with
$\mathfrak m_{p_0}$. This recovers the displayed $\mathcal L$
with its original cusp condition. These transformation and order
calculations are the ones used to construct the retained $\mu$;
no topological conclusion about $X_c$ is involved.

For completeness, exactness on the finite stalks in (PS24) retains
the off-diagonal entry of $R_g$. Given an equivariant first component
$u(z)$, begin with the vector $(u(z),0)^{\mathsf t}$ and average it
under the finite group action on vector-valued holomorphic functions.
The resulting vector is equivariant, and its first component is
still $u$. Thus projection is surjective there. Its kernel is
exactly the invariant second-component functions, which are
holomorphic functions of $s_j^{m_j}=t_j$, multiplied by $e$.
The ordinary and cusp stalks are immediate. This also proves that
$\mathscr V$ is locally free of rank two: lift a local frame of
$\mathcal L$ and join it to $e$.

Both $H^1(B,\mathcal O_B)$ and
$H^1(B,\mathcal O_B(-p_0))$ vanish. One can compute these two
groups without a degree approximation: on the two standard affine
charts of $\mathbb P^1$, split a holomorphic Laurent series on
their overlap into its nonnegative and negative powers. For
$\mathcal O_B$ these are the two coboundary contributions, and
for $\mathcal O_B(-p_0)$ its transition factor shifts the second
contribution by one, still leaving no missing exponent. The positive
part converges on the full first affine chart and the transformed
negative part on the full second affine chart. This proves the two
vanishings, equivalently the usual additive Cousin calculation on
these charts. The exact sequence (PS24) now gives
\[
                             H^1(B,\mathscr V)=0.
 \tag{PS25}\label{wb:pshift:eq:no-analytic-obstruction}
\]

Here is how to choose the correction entering the maps. Choose
local sections $h_\alpha$ of $\mathscr H_n$, taking the displayed
$a_j,a_0$ on sufficiently small special discs and affine charts
elsewhere. Their differences $d_{\alpha\beta}=h_\beta-h_\alpha$
are a Cech $1$-cocycle with values in $\mathscr V$. Project it to
$\mathcal L$ and solve that additive cocycle, using its just proved
$H^1$ vanishing. Lift the resulting local $\mathcal L$ sections to
$\mathscr V$ by the local lifting construction in (PS24). Subtracting
their coboundary from $d$ leaves a scalar $\mathcal O_Be$ cocycle.
Solve that scalar additive cocycle by its $H^1$ vanishing. This
produces sections $v_\alpha$ with
$d_{\alpha\beta}=v_\beta-v_\alpha$. The local sections
$h_\alpha-v_\alpha$ agree and define a global section
$a\in\mathscr H_n(B)$. This is a two-stage construction of the
required correction, not an unproved additional gluing condition.
It gives holomorphic vectors $\eta_j,\eta_0$ with
\[
 a=a_j+\eta_j\quad(j=1,2),\qquad a=a_0+\eta_0
                              \quad\hbox{at the cusp}.
 \tag{PS26}\label{wb:pshift:eq:global-correction}
\]
Any choices made in the additive solutions give maps with the same
asserted properties. Their differences as sections of this affine
bundle lie in $H^0(B,\mathscr V)=\mathbb C e$: this follows from
(PS24), $H^0(B,\mathcal L)=0$, and $H^0(B,\mathcal O_B)=\mathbb C$.

Projecting $a$ modulo $\Pi_c\Lambda$ gives a holomorphic section of
the middle torus family because every $\ell_g$ in (PS23) is integral.
The promised three kinds of local maps are therefore
\[
\begin{aligned}
 F_{\mathrm{mid}}[z,\zeta]&=[z,\zeta+a(z)],\\
 F_j[z,\zeta]&=[z,\zeta+b_j(z)+\eta_j(z)]\quad(j=1,2),\\
 F_0[x,t_c]&=[\exp(2\pi i\eta_0(t_c))t_c^{q_n}x,t_c].
\end{aligned}
\tag{PS27}\label{wb:pshift:eq:full-isomorphism}
\]
The first descends by (PS23) and equality of the integer-shift
lattices. The finite maps descend by (PS10), (PS14), and the exact
covariance of $\eta_j$. The cusp map descends by (PS19); its
holomorphic-unit factor extends by torus multiplication on every
original chart, and its integer factor extends by the fan map in
(PS22). Each map is a biholomorphism, with inverse given by the
negative translation or the inverse torus multiplier and fan map.

On a finite overlap, write the original gluing as
$h_{j,c}(\zeta)=\zeta+\sigma_{j,c}$. Then (PS17), (PS26), and
(PS27) give the equality of actual torus maps
\[
                     h_{j,c+n}F_j=F_{\mathrm{mid}}h_{j,c}.
 \tag{PS28}\label{wb:pshift:eq:finite-overlap-identity}
\]
Both sides add
$\sigma_{j,c}+a=\sigma_{j,c+n}+b_j+\eta_j$.
On the cusp overlap the original map is exponentiation of $\zeta$;
(PS22) and (PS26) give
$\exp(2\pi ia)=t_c^{q_n}\exp(2\pi i\eta_0)$, exactly the
third line of (PS27). Thus that overlap also agrees, including
all nonconstant period units. The maps and their inverses glue to
a global biholomorphism $X_c\to X_{c+n}$ over $B$.

The choices can be made coherently for every parameter and every even
shift. Both initial cocycle columns in (PS13) have $\gamma=0$.
Because every $A_g$ preserves $\gamma$, (PS20) shows that
$\gamma(\ell_g)=0$ for every word. Therefore the affine term
$\Pi_c(gz)\ell_g$ in (PS23) is independent of $c$. The displayed
$a_j,a_0$, and the vector sheaf $\mathscr V$ are also independent
of $c$. Perform the additive construction once for $n=2$, keeping
one resulting section $a^{(2)}$. For $n=2k$ choose
$a^{(2k)}=k a^{(2)}$ and multiply its local correcting vectors by
$k$. All their affine cocycles and local formulas are then exactly
those for $n=2k$. The resulting maps satisfy the full group law
\[
 F_{2h,c+2k}\circ F_{2k,c}=F_{2(h+k),c},\qquad
                 F_{0,c}=\mathrm{id}_{X_c}.
 \tag{PS28a}\label{wb:pshift:eq:coherent-shift-action}
\]
This follows by adding the actual translation vectors in the first
two lines of (PS27), and by multiplying their units and adding the
integer exponents in its third line. The formulas are jointly
holomorphic in $c$, since their corrections are independent of it.

\subsubsection{Intrinsic necessity and the resulting parameter space}

The analytically computed complete linear system
$|\omega_{X_c}^{-2}|$ is the original map $f_c$, by
\cref{wb:thm:canonical-ring-recovery}. A biholomorphism preserves its
canonical bundle, the complete space of sections, and its evaluation
map. Consequently it induces $\alpha\in\operatorname{PGL}_2(\mathbb C)$
with $f_{c'}F=\alpha f_c$. The three special fibres have distinct
intrinsic types: the cusp is reduced and nonnormal, while the other
two have multiplicities three and four. All remaining fibres are
smooth reduced tori. Thus $\alpha$ fixes $p_0,p_1,p_2$ individually.
In the original coordinate, a fractional-linear map fixing
$\infty,0,1$ has zero lower-left and upper-right entries and equal
diagonal entries; hence $\alpha=\mathrm{id}_B$.

Now \cref{wb:pshift:lem:centralizer} forces $c'-c=n\in\mathbb Z$
and linear part $\pm I_2$. Equations (PS10)--(PS11), obtained on
the original normalized finite covers, exclude $-I_2$ and then
exclude odd $n$ for $+I_2$. Conversely (PS27)--(PS28) construct
the global map for every even $n$. This proves
\cref{wb:pshift:thm:classification} without converting a failed
particular map into a nonisomorphism assertion.

As a complex parameter space for the isomorphism classes represented
by this retained family, the quotient $\mathcal P/(2\mathbb Z)$ has
the explicit coordinate
\[
 w=\exp(-\pi ic),\qquad 0<|w|<\exp(-\pi M).
 \tag{PS29}\label{wb:pshift:eq:parameter-disc}
\]
Indeed $|w|=\exp(\pi\operatorname{Im}c)$, and two complex numbers
have the same value of this exponential exactly when they differ
by an even integer. Every point in the displayed punctured disc
has a logarithm giving a $c\in\mathcal P$. The exponential has
nonzero derivative, so these local inverse logarithms give the
claimed quotient complex structure. The integer-period quotient
seen by the middle family instead has coordinate $w^2$; it is
covered two-to-one by (PS29). Its two lifts are distinguished by
the exact order-two summand in (PS12). This statement classifies
the parameters of the retained family; it does not assert that
the family exhausts all deformations or all complex threefolds
with the same singular-fibre types.



\subsection{Re-derivation of Theorem~\ref{fam:src:A.5.1}}\label{app:rederive:fam-src-A.5.1}
\begin{proof}
\ref{fam:src:A.1.1} makes the base intrinsic. \ref{fam:src:A.2.1} gives the necessary parity and excludes the negative linear part. \ref{fam:src:A.3.1} supplies the inverse maps, all collar identities, and descent of the jointly holomorphic family. \ref{fam:src:A.4.2} proves its infinitesimal effectiveness. Finally $|e^{-\pi ic}|=e^{\pi\operatorname{Im}c}$ and equality of exponentials is exactly difference in $2\mathbb Z$, so the coordinate has precisely the stated image and fibres.
\end{proof}

\subsection{Re-derivation of Theorem~\ref{fam:src:A.5.2}}\label{app:rederive:fam-src-A.5.2}
\begin{proof}
We give the calculation, including the singular-fibre maps. The
corresponding calculation appears in the version 1
\cite[Section 6, especially Proposition 6.10, Lemma 6.11,
Theorem 6.12 and Proposition 6.16]{DuEtAl2026v1}; those deformation
statements are omitted from version two. We use the verified maps and
cohomology calculation below, rather than infer their validity from the
earlier theorem label.

\emph{The base and its line bundles.}
Write $X=X_c$ and form the analytic root curve
\[
 \mathfrak B=\sqrt[3]{(B,p_1)}\mathbin{\times_B}\sqrt[4]{(B,p_2)},
 \qquad \rho:\mathfrak B\longrightarrow B.
\]
A root object is a line bundle with section whose indicated power is the
pullback of the section cutting out $p_j$. The actual equalities
$f^*p_1=3S_1$ and $f^*p_2=4S_2$ supply these objects as
$(\mathcal O_X(S_j),1_{S_j})$, and hence supply a morphism
$\bar f:X\to\mathfrak B$ with $\rho\bar f=f$.
The explicit weighted quotient
\[
 \mathfrak B\simeq[(\mathbb C^2\setminus\{0\})/\mathbb C^*],
 \qquad \lambda(x,y)=(\lambda^3x,\lambda^4y),
 \qquad t=\frac{y^3}{y^3-x^4},
\]
identifies the roots with $(\mathcal O(4),y)$ and
$(\mathcal O(3),x)$ and the ordinary point $p_0$ with
$y^3-x^4=0$. In particular $\rho^*\mathcal O_B(1)=\mathcal O(12)$.
This isomorphism follows on each root chart by the displayed root
sections; away from their zeros both sides have the same ordinary coarse
coordinate, so the chart isomorphisms glue.

For completeness the analytic line-bundle computation can be made on
the coarse curve. Removing the two stabilizer characters from any line
bundle makes it descend to $B$: on each root chart, averaging a local
frame whose central value is nonzero gives an invariant frame still
nonzero at the centre, and its invariant coefficients are holomorphic
in $s_j^{m_j}$. These frames give the descended line and its transition
functions. Consequently its Picard group has
generators $H,R_1,R_2$ and relations $3R_1=H=4R_2$.
Thus it is freely generated by $R_1-R_2$, denoted $\mathcal O(1)$;
$H=12\mathcal O(1)$, $R_1=4\mathcal O(1)$ and
$R_2=3\mathcal O(1)$. The chart descent just used is exact because
invariants of a finite group over $\mathbb C$ are an exact functor.
With the clockwise root generators, a local invariant section of
$\mathcal O(k)$ has vanishing exponents
\[
 r_3(k)\in\{0,1,2\},\quad r_3(k)\equiv k\pmod3,
 \qquad
 r_4(k)\in\{0,1,2,3\},\quad r_4(k)\equiv-k\pmod4.
\]
Dividing by those local powers gives the precise coarse direct image
\begin{equation}\label{fam:eq:ab-root-degree}
 \rho_*\mathcal O(k)=\mathcal O_B\bigl(d(k)\bigr),\qquad
 d(k)=\frac{k}{12}-\frac{r_3(k)}3-\frac{r_4(k)}4.
\end{equation}
Indeed $k-4r_3-3r_4=12d(k)$ identifies the line bundle with
$\rho^*\mathcal O_B(d(k))\otimes R_1^{r_3}\otimes R_2^{r_4}$.
The canonical root sections generate the invariant direct image of
the last two factors; division by their local powers is the displayed
line-bundle isomorphism on the coarse charts. Finite invariants also give
$R^q\rho_*=0$ for $q>0$, as seen on each Stein quotient chart.
For $k=0$ the degree is zero; for
$k=1,-1,-2,-5,-6$ it is $-1$. It follows directly from the
cohomology of $\mathbb P^1$ that these five line bundles are acyclic,
whereas $H^0(\mathcal O)=\mathbb C$ and $H^1(\mathcal O)=0$.
The local equation $dt_j=m_js_j^{m_j-1}ds_j$ gives
\[
 K_{\mathfrak B}=\mathcal O(-24+2\cdot4+3\cdot3)
 =\mathcal O(-7),\qquad
 T_{\mathfrak B}(-p_0)=\mathcal O(-5).
\]

\emph{The actual base changes and cusp cohomology.}
On a finite root disc $\Delta_{s_j}$ the local presentations are
$[J'_j/\mu_{m_j}]\to[\Delta_{s_j}/\mu_{m_j}]$.
The fibre-product identity
\[
 [J'_j/\mu_{m_j}]
 \mathbin{\times_{[\Delta_{s_j}/\mu_{m_j}]}}\Delta_{s_j}
 \simeq J'_j
\]
is induced by the quotient torsor and its equivariant map to the disc;
its inverse sends a torsor with a chosen root-chart trivialization to
the corresponding point of $J'_j$. The action on $J'_j$ is free by
\ref{fam:src:A.1.2}. Thus this is an etale atlas base change and is exactly the
smooth torus family, with its given deck generator. Over the cusp the
base change is the original cusp family. Locally it is
$q=z_0z_1z_2$; the homomorphism
$\mathbb C\{q\}\to\mathbb C\{z_0,z_1,z_2\}$ is flat because its
target is torsion-free over the discrete valuation ring.
These charts prove that $\bar f$ is representable, proper and flat,
with connected Gorenstein fibres of dimension two. Properness on the
cusp is the compact representative construction of \ref{fam:src:A.3.1}, and on the
finite charts it is the compact torus quotient. Also
$\bar f^*p_0=W$ as reduced Cartier divisors.

Let $\nu:E\to W$ be the normalization, $E=\mathrm{dP}_6$, let
$C_h,C_v,C_d$ be the three double-curve normalizations, and let
$P_+,P_-$ be the two triple points. Each $C_e$ is $\mathbb P^1$.
Write $\rho_e^0,\rho_e^1:C_e\to E$ for its two paired boundary maps
and $j_e=\nu\rho_e^0=\nu\rho_e^1$. The sheaf sequence is
\begin{equation}\label{fam:eq:ab-cusp-resolution}
 0\longrightarrow\mathcal O_W\longrightarrow\nu_*\mathcal O_E
 \xrightarrow{\delta^0}\bigoplus_{e=h,v,d}j_{e*}\mathcal O_{C_e}
 \xrightarrow{\delta^1}\mathbb C_{P_+}\oplus\mathbb C_{P_-}
 \longrightarrow0,
\end{equation}
where $\delta^0(s)_e=(\rho_e^1)^*s-(\rho_e^0)^*s$ and
$\delta^1$ evaluates $t_h-t_v+t_d$ at the two ordered triple points.
Its exactness is local: for $xyz=0$ it is the alternating restriction
complex of the three coordinate planes, their pairwise lines and their
common point. Decomposing convergent power series by the variables
occurring in each monomial gives the augmented simplex complex, which
is exact. The one- and two-branch cases are the same calculation with
fewer planes. This verifies exactness for the actual boundary maps,
including the torus multipliers in their gluing.

All positive-degree structure-sheaf cohomology of $E,C_e$ vanishes.
Consequently \eqref{fam:eq:ab-cusp-resolution} computes cohomology by
\[
 \mathbb C\xrightarrow{0}\mathbb C^3
 \xrightarrow{\left(\begin{smallmatrix}1&-1&1\\1&-1&1\end{smallmatrix}\right)}
 \mathbb C^2.
\]
It gives $(h^0,h^1,h^2)(\mathcal O_W)=(1,2,1)$.
We also need its multiplication. The branch-intersection resolution
is the augmented semisimplicial sheaf algebra with its restriction face
maps; its local exactness makes the augmentation a multiplicative
quasi-isomorphism. Its Alexander--Whitney product therefore computes
the sheaf-cohomology cup product. In the ordered triangles
$[0,e_1,e_2]$ and $[e_1,e_2,e_1+e_2]$ it is
\[
 \alpha\smile\beta=(\alpha_h\beta_d,\alpha_d\beta_h),
 \qquad [\alpha]\smile[\beta]
 =\alpha_h\beta_d-\alpha_d\beta_h
 \quad\hbox{in }\mathbb C^2/\mathbb C(1,1).
\]
The cocycles $(1,1,0)$ and $(0,1,1)$ pair to one. Thus
$\bigwedge^2H^1(\mathcal O_W)\to H^2(\mathcal O_W)$ is an
isomorphism, including its product rather than just its dimensions.

\emph{Vertical fields and the direct images.}
Set $\mathcal T=T_X(-\log W)$ and
$\mathcal V=\ker(d\bar f^{\log})$.
The vertical transitions on the core are precisely $R_g(z)$.
Finite collar translations have vertical differential the identity,
and at the cusp the frame $\partial_{\zeta_i}$ corresponds to
$2\pi i x_i\partial_{x_i}$. The cusp transformations
$x\mapsto u(q)q^{B_0\lambda}x$ fix the coefficients of these
relative logarithmic fields. These exact identifications construct
a bundle $\mathcal E$ on $\mathfrak B$ with
$\mathcal V=\bar f^*\mathcal E$.
The fixed vector $e=(0,1)^t$ is nonzero in every logarithmic frame,
so there is an exact sequence
\[
 0\longrightarrow\mathcal O\xrightarrow{e}\mathcal E
 \longrightarrow L\longrightarrow0.
\]
The factors of $L$ are $-1/\tau$ and $1/\tau$. Therefore the usual
weight-twelve modular discriminant makes $1/\Delta(\tau)$ a
meromorphic section of $L^{12}$. It has no finite zero or pole,
and its cusp pole is simple: $\Delta(\tau)=q\,u(q)$ with
$u(0)\ne0$. Hence $L^{12}=\mathcal O(-p_0)=\mathcal O(-12)$.
The Picard calculation gives $L=\mathcal O(-1)$, and the extension
splits since $H^1(\mathcal O(1))=0$. Thus
\[
 \mathcal E\simeq\mathcal O\oplus\mathcal O(-1).
\]
There is also the exact sequence
\begin{equation}\label{fam:eq:ab-log-tangent}
 0\longrightarrow\bar f^*\mathcal E\longrightarrow\mathcal T
 \longrightarrow\bar f^*\mathcal O(-5)\longrightarrow0.
\end{equation}
On the smooth root charts it is the relative tangent sequence. At the
cusp its last map sends $\sum a_i z_i\partial_{z_i}$ to
$(\sum a_i)q\partial_q$, so it is onto and its kernel is the
two height-zero lattice directions already used above. Taking its
determinant and using $W=\bar f^*p_0$ gives the actual relative
dualizing line
\begin{equation}\label{fam:eq:ab-relative-dualizing}
 \omega_{X/\mathfrak B}=\bar f^*(\det\mathcal E)^{-1}.
\end{equation}
In particular this computation gives
$K_X=\bar f^*\mathcal O(-6)=\mathcal O_X(-2S_2)$, with the last
identification supplied by $\bar f^*R_2=\mathcal O_X(S_2)$;
it agrees through this actual pullback with \ref{fam:src:A.1.1}.

The fibres of $\bar f$ all have structure-sheaf cohomology dimensions
$(1,2,1)$: on ordinary and finite root charts they are tori, and at
the cusp we have just computed them. Analytic proper flat base change
on these complex-space charts gives locally free $R^q\bar f_*\mathcal O_X$
with these fibres. Its canonical comparison maps agree on the atlas
double overlaps, so they descend to $\mathfrak B$. The unit gives
$R^0\bar f_*\mathcal O_X=\mathcal O$.
The source charts are smooth and the fibres are Cartier divisors,
so relative analytic duality applies in its Gorenstein form. The
relative trace pairing, compatible on the same overlaps, and
\eqref{fam:eq:ab-relative-dualizing} give
\[
 (R^2\bar f_*\mathcal O_X)^\vee
 =\bar f_*\omega_{X/\mathfrak B}
 =(\det\mathcal E)^{-1},
 \qquad R^2\bar f_*\mathcal O_X=\mathcal O(-1).
\]
Here local freeness removes the higher sheaf-Ext terms in the duality
spectral sequence. This is precisely the pointwise base-change and
top-degree duality application spelled out in the cited Lemmas 6.8--6.9.

Put $G=R^1\bar f_*\mathcal O_X$. Cup product and the verified fibre
comparisons show $\det G=R^2\bar f_*\mathcal O_X=\mathcal O(-1)$;
at $W$ this uses the nonzero pairing above.
We construct the required subline of $G$. On an ordinary fibre,
constant lattice cocycles identify
\[
 H^1(F_t,\mathcal O_{F_t})
 =\Lambda_{\mathbb C}^*/\operatorname{rowspan}_{\mathbb C}\Pi_c.
\]
Indeed a coboundary among constant cocycles is given by an entire
function with constant period differences. Its derivatives are
periodic holomorphic functions on the compact torus, hence constant;
the function is affine linear. This gives exactly the displayed
kernel, and both sides then have dimension two. The class of
$\gamma=(1,0,0,0)$ is nonzero, since evaluation on columns
$\widehat w,\widehat\delta$ would force both coefficients of a
putative period-row combination to vanish.
The matrices satisfy $\gamma A_j=\gamma$. Translations act trivially
on $H^1(\mathcal O)$, so this holomorphic class extends on both finite
root charts and respects their gluing.

At the cusp $\gamma$ kills
$\Lambda_{\mathrm{tor}}=\mathbb Z\langle\widehat w,\widehat\delta\rangle$.
For local sections of the free toric quotient with changes
$s_i=\Psi_{\lambda_{ij}}s_j$, the constants
$\gamma(\lambda_{ij})$ therefore form a holomorphic Cech cocycle
on the entire neighbourhood. With
$B_0=\left(\begin{smallmatrix}0&1\\-1&0\end{smallmatrix}\right)$,
the edges $e_1,e_2,e_2-e_1$ have deck vectors
$(0,1),(-1,0),(-1,-1)$. The central cocycle values are consequently
$(0,-1,-1)$, which pair to one with $(1,0,-1)$.
Its value on $W$ is nonzero. On the punctured overlap the exponential
map kills precisely $\Lambda_{\mathrm{tor}}$, so its full-lattice
cocycle is exactly $\lambda\mapsto\gamma(\lambda)$, proving equality
with the preceding class. Over the Stein cusp disc, the Leray edge
map $H^1(\bar f^{-1}D_0,\mathcal O_X)\to H^0(D_0,G)$ is an
isomorphism, since $H^1(D_0,\mathcal O)=H^2(D_0,\mathcal O)=0$.
Thus this is a holomorphic, nonvanishing extension of $[\gamma]$.
It supplies a global subline and exact sequence
\[
 0\longrightarrow\mathcal O\xrightarrow{[\gamma]}G
 \longrightarrow\mathcal O(-1)\longrightarrow0.
\]
Again $H^1(\mathcal O(1))=0$ splits it. We have proved
\begin{equation}\label{fam:eq:ab-direct-images}
 R^q\bar f_*\mathcal O_X=
 \begin{cases}
 \mathcal O,&q=0,\\
 \mathcal O\oplus\mathcal O(-1),&q=1,\\
 \mathcal O(-1),&q=2.
 \end{cases}
\end{equation}

\emph{From logarithmic to ordinary tangent cohomology.}
The projection formula gives, for $q=0,1,2$ respectively,
\[
 R^q\bar f_*\mathcal V:
 \quad\mathcal O\oplus\mathcal O(-1),\quad
 \mathcal O\oplus2\mathcal O(-1)\oplus\mathcal O(-2),\quad
 \mathcal O(-1)\oplus\mathcal O(-2).
\]
For $\bar f^*\mathcal O(-5)$ the three expressions are
$\mathcal O(-5)$, $\mathcal O(-5)\oplus\mathcal O(-6)$ and
$\mathcal O(-6)$. Formula \eqref{fam:eq:ab-root-degree} makes every
nontrivial summand in this list acyclic. Leray, whose base has
cohomological dimension one, therefore gives
$h^\bullet(\mathcal V)=(1,1,0,0)$ and zero cohomology for the
horizontal term. Sequence \eqref{fam:eq:ab-log-tangent} yields
$h^\bullet(\mathcal T)=(1,1,0,0)$.

The normalization map $\nu:E\to X$ is an immersion on each local
branch. Let $N_\nu=\operatorname{coker}(T_E\to\nu^*T_X)$.
The exact normal-derivative map is
\begin{equation}\label{fam:eq:ab-normal-quotient}
 0\longrightarrow\mathcal T\longrightarrow T_X
 \longrightarrow\nu_*N_\nu\longrightarrow0.
\end{equation}
For a local equation $z_0\cdots z_{r-1}=0$, it sends
$\sum b_i\partial_{z_i}$ to the $r$ independent normal coefficients
$b_i|_{z_i=0}$. These can be lifted independently, proving
surjectivity; the kernel is exactly the logarithmic tangent module.
This also specifies the global map under changes of branch charts.
The cusp three-form
\[
 \Omega=d\log x_1\wedge d\log x_2\wedge dq
\]
is invariant under $x\mapsto u(q)q^{B_0\lambda}x$ because the
extra logarithmic differentials are multiples of $dq$. In any
unimodular maximal cone the height-one equation is $q=z_0z_1z_2$,
and the determinant-one character change gives
$\Omega=\pm dz_0\wedge dz_1\wedge dz_2$.
It therefore trivializes $K_X$ on the entire cusp neighbourhood,
including every branch intersection. Taking determinants of the
immersion tangent sequence now gives the explicit line identification
$N_\nu=K_E\otimes(\nu^*K_X)^{-1}\simeq K_E$.
Since $E$ is the rational projective surface $\mathrm{dP}_6$,
Serre duality gives $h^\bullet(E,K_E)=(0,0,1)$.
The cohomology sequence of \eqref{fam:eq:ab-normal-quotient} proves
$h^\bullet(X,T_X)=(1,1,1,0)$, as asserted.

\emph{The reduced deformation germ.}
\ref{fam:src:A.4.2} gives a nonzero Kodaira--Spencer vector at this same $c$;
the just computed $h^1=1$ makes its map an isomorphism.
Embed the actual Kuranishi germ in its one-dimensional tangent space
and write its local analytic algebra as $\mathbb C\{t\}/I$.
The existing proper family has a local classifying map, whose pullback
sends $t$ to $g(s)$ with $g'(0)\ne0$. The holomorphic inverse-function
theorem makes $\mathbb C\{t\}\to\mathbb C\{s\}$, $t\mapsto g(s)$,
an isomorphism. Every element of $I$ maps to zero, so $I=0$.
This proves smoothness and reducedness even though $h^2(T_X)=1$.
The classifying map is consequently a local isomorphism of germs;
completeness of the Kuranishi deformation then gives the claimed
miniversality and exhaustion of sufficiently small deformations.
\end{proof}

\subsection{Re-derivation of Proposition~\ref{fam:src:B.4.1}}\label{app:rederive:fam-src-B.4.1}
\begin{proof}
\ref{fam:src:A.1.1} and \ref{fam:src:A.2.1} give base identity, $n=0$ and fibre-linear part $I$. Let $\ell_1,\ell_2$ be the integral defects of a lift of its translating section to the universal cover of $B^\circ$. Finite extension gives $N_j\ell_j=0$. The saturated image calculations in \ref{fam:src:A.1.2} give the integral equality $\ker N_j=(I-A_j)\Lambda$. The cusp valuation argument (B2.3) gives $\ell_0\in\Lambda_{\rm tor}$. Write
\[
 \ell_1=ae_1+bf_1,\qquad \ell_2=ce_2+df_2.
\]
The $\widehat u$ coordinate of $\ell_1+A_1\ell_2$ is $a-c-2d$, so the cusp condition is exactly $a=c+2d$. Changing the global lift by a period $\lambda$ changes this pair by the simultaneous coboundary $((I-A_1)\lambda,(I-A_2)\lambda)$. In the four displayed coefficient coordinates its matrix is
\[
 C=\begin{pmatrix}-6&1&-1&0\\2&-1&0&0\\0&1&1&0\\-3&0&-1&0\end{pmatrix}.
\]
The lattice of all allowed pairs has basis matrix
$K=\left(\begin{smallmatrix}0&1&2\\1&0&0\\0&1&0\\0&0&1\end{smallmatrix}\right)$. The first three columns of $C$ expressed in this basis have determinant one. Hence every allowed pair is an integral coboundary; the computation is exact, including its saturation, and is independently checked by Hermite form in \texttt{ab\_exact.py}. After changing the lift, all its defects vanish.

The corrected lift extends as an equivariant holomorphic vector on the finite discs. A holomorphic local lift exists on the whole root disc by lifting the translating section through its covering vector bundle. Its defect is in $\ker N_j=(I-A_j)\Lambda$; subtracting an integral period kills that defect while preserving holomorphy. The corrected global lift and this local defect-zero lift differ by an invariant period, which is holomorphic through the centre. At the cusp, the local toric logarithm in (B2.3) is $sq+\eta(t_c)$. The chosen global lift differs from it by $\Pi_c(s)\lambda=sB_0\bar\lambda+C(t_c)\bar\lambda+\lambda_{\rm tor}$. Defect zero gives exactly $q+B_0\bar\lambda=0$, so the logarithmic terms cancel and the remaining vector is holomorphic. This proves the extension at the cusp as well. We therefore obtain a global section of $\mathcal V$. Its exact sequence in \ref{fam:src:A.3.1} gives $H^0(B,\mathcal V)=\mathbb Ce$. Conversely each such section defines the displayed translations on the finite covers and torus multiplication at the cusp, hence a global automorphism. Its kernel consists of global invariant periods. The common fixed lattice is $\mathbb Zd$, and $\Pi_cd=e$, giving kernel exactly $\mathbb Z$. This proves the full group assertion and its connectedness.
\end{proof}


\section{Tangent cohomology and deformations}

We retain the standard filling data $(\ell _0,\ell _1,\ell _2)=(0,1,-1)$
and an arbitrary admissible $c$. All sheaves in this section are analytic,
and all cohomology groups have coefficients in $\C$. The integral twists
below mean the actual line bundles
\[
 M_n=\OO_X\bigl(n(S_1-S_2)\bigr),\qquad n\in\Z.
\]
We set $\mathfrak g=g_\gamma$. Thus, in the positively normalized
exponential coordinate of \ref{inv:C3.2},
$M_n=\exp(n\mathfrak g)$. The computations use the root base to retain the
nonreduced fibres. Passing to its coarse space is a separate, exact
invariant operation; it is never identified with taking the dual of an
invariant lattice.

\begin{proposition}[The root base and its explicit cohomology]\label{inv:D1.1}
\StatementRecord{inv:D1.1}{proved}{Historical events: GPT-6 Astra (OpenAI), ChatGPT Pro mode, effort Pro, S6\_09\_PAPERS, 17 November 2026: targeted proof check. Also checked in this writing. Pre-continuation writing events: GPT-6 Astra (OpenAI), ChatGPT Pro mode, effort Pro, S6\_10\_PAPERS, 18 November 2026: targeted proof check. Continuation metadata events: GPT-6 Astra (OpenAI), ChatGPT Pro mode, effort Pro, S6\_10\_PAPERS continuation, 19 November 2026: used and rechecked in this run; status-vocabulary reconciliation; metadata reconciliation only; no new mathematical check. \emph{Pre-continuation recorded proof scope.} Read full root-base, logarithmic tangent, direct-image and twisted-Dolbeault proofs. Re-derived coarse degree function, local normalization quotient, resonance equation, Fourier contraction comparison, joint order-three coefficient 123ik/(64pi), stated local modules, and jump-locus finiteness. Exact Koszul and twist-program computations retain their separately recorded replay scopes.}{the hypotheses, notation and definitions stated here; the explicit constructions and calculations in the proof}

Let $\rho:\mathcal B\to B$ be the analytic root curve with orders $3$ and
$4$ at $p_1,p_2$. Write $R_1,R_2$ for its tautological root line bundles.
Its Picard group is infinite cyclic, with generator
$R=R_1\otimes R_2^{-1}$, and
\[
 R_1=R^4,\quad R_2=R^3,\quad \rho^*\OO_B(1)=R^{12},
 \quad K_{\mathcal B}=R^{-7}.
\]
We abbreviate $R^n$ to $\OO_{\mathcal B}(n)$. Choose sections
$x\in H^0(R^3)$ and $y\in H^0(R^4)$ vanishing at $p_2$ and $p_1$,
respectively. For every integer $n$, put
\[
 r_1(n)\equiv n\pmod3,\quad 0\le r_1(n)<3,
 \qquad r_2(n)\equiv3n\pmod4,\quad0\le r_2(n)<4,
\]
\[
 d(n)=\frac{n-4r_1(n)-3r_2(n)}{12}.
\]
Then
\[
 \rho_*\OO_{\mathcal B}(n)=\OO_B(d(n)),\qquad
 H^i(\mathcal B,\OO(n))=H^i(B,\OO(d(n))).                 \tag{D.1}
\]
Equivalently, $h^0(\OO(n))$ is the number of pairs
$(a,b)\in\Z_{\ge0}^2$ with $3a+4b=n$, and
$h^1(\OO(n))=h^0(\OO(-n-7))$.
\end{proposition}



The root-curve Picard group and graded cohomology are
\citet[Lemma~6.1]{DuEtAl2026v1}. We re-derive them and retain the explicit
coarse invariant-descent function $d(n)$ in the same calculation.

\begin{proof}
\begin{enumerate}[label=\arabic*.,leftmargin=*]
\item On the root charts the group acts on a disc coordinate by a primitive
$m_j$th root of unity. The stabilizer character of any line bundle is
removed uniquely by tensoring with $R_1^{-r_1}R_2^{-r_2}$, with the
exponents in the displayed ranges. A line bundle with trivial stabilizer
characters descends: an equivariant nonzero local frame can be averaged,
and its invariant frame gives the descent transition function. It is
therefore the pullback of a unique line bundle on $B$. Consequently
\[
 \operatorname{Pic}(\mathcal B)=
 \langle H,R_1,R_2\mid 3R_1=H,\;4R_2=H\rangle=\Z R.
\]
Here and below additive notation for this presentation means tensor
product. The determinant of the coordinate change $t=s^m$ gives
$K_{\mathcal B}=\rho^*K_B\otimes R_1^2R_2^3=R^{-24+8+9}$.

\item An invariant section of $R_j^{r_j}$ on its chart is $s^{r_j}$ times a
holomorphic function of $t=s^{m_j}$. It follows, by the local invariant
comparison on overlaps, that
$\rho_*(\rho^*\OO(k)R_1^{r_1}R_2^{r_2})=\OO(k)$.
In characteristic zero taking finite-group invariants is exact: the
averaging idempotent splits the invariant submodule. Thus there are no
higher direct images for $\rho$, and (D.1) follows. This proves the
cohomology comparison in the direction used here, not merely equality
of Euler characteristics. Finally every monomial of degree $n$ is
$x^{r_2}y^{r_1}$ times a homogeneous polynomial of degree $d(n)$ in
$x^4,y^3$. This gives the stated count of sections. The negative-degree
Cech complex on $x\ne0$, $y\ne0$ has cokernel represented by
$x^{-a}y^{-b}$ with $a,b\ge1$ and $3a+4b=-n$; replacing $(a,b)$ by
$(a-1,b-1)$ proves the $h^1$ formula without an additional duality
assumption. The two calculations agree exactly for every integer $n$;
the program \texttt{code/d\_exact.py} also compares them for
$-240\le n\le240$.
\end{enumerate}
\end{proof}

\begin{proposition}[Vertical and logarithmic tangent bundles]\label{inv:D1.2}
\StatementRecord{inv:D1.2}{proved}{Historical events: GPT-6 Astra (OpenAI), ChatGPT Pro mode, effort Pro, S6\_09\_PAPERS, 17 November 2026: targeted proof check. Also checked in this writing. Pre-continuation writing events: GPT-6 Astra (OpenAI), ChatGPT Pro mode, effort Pro, S6\_10\_PAPERS, 18 November 2026: targeted proof check; used and rechecked in this run. Continuation metadata events: GPT-6 Astra (OpenAI), ChatGPT Pro mode, effort Pro, S6\_10\_PAPERS continuation, 19 November 2026: used and rechecked in this run; status-vocabulary reconciliation; metadata reconciliation only; no new mathematical check. \emph{Pre-continuation recorded proof scope.} Read full root-base, logarithmic tangent, direct-image and twisted-Dolbeault proofs. Re-derived coarse degree function, local normalization quotient, resonance equation, Fourier contraction comparison, joint order-three coefficient 123ik/(64pi), stated local modules, and jump-locus finiteness. Exact Koszul and twist-program computations retain their separately recorded replay scopes. Matched the previously successful unchanged d\_exact.py and d\_koszul.py replays to this full tangent/Koszul proof:481 weighted cohomology counts, root-character invariants, polynomial module relations and the joint coefficient123i/(64pi). The earlier event saying replay pending is superseded at exactly this finite-algebra scope; its complete analytic comparison review remains separate.}{the hypotheses, notation and definitions stated here; \cref{inv:D1.1}}

There is a canonical representable proper flat holomorphic map
$\bar f:X\to\mathcal B$ with $f=\rho\bar f$. If
\[
 \mathcal T=T_X(-\log W),\qquad
 \mathcal V=\ker\bigl(\mathcal T\to
                   \bar f^*T_{\mathcal B}(-p_0)\bigr),
\]
then, for rank-two vector bundles $E,G$ on $\mathcal B$,
\[
 \mathcal V=\bar f^*E,\quad E=\OO\oplus\OO(-1),\qquad
 G=R^1\bar f_*\OO_X=\OO\oplus\OO(-1),                 \tag{D.2}
\]
\[
 R^0\bar f_*\OO_X=\OO,\qquad R^2\bar f_*\OO_X=\OO(-1),
\]
and the following is an exact sequence of locally free sheaves on $X$:
\[
 0\longrightarrow\bar f^*E\longrightarrow\mathcal T
 \longrightarrow\bar f^*\OO(-5)\longrightarrow0.       \tag{D.3}
\]
The trivial summands in $E,G$ are generated by the vertical field $\xi$
and the constant integral cocycle $\gamma$, respectively.
\end{proposition}



The root lift, vertical splitting, logarithmic quotient and
structure direct images are re-derived after reading
\citet[Lemmas~6.2 and 6.7; Propositions~6.6 and 6.10]{DuEtAl2026v1}.
The complete root and ordinary-sheaf arguments are retained below.

\begin{proof}
\begin{enumerate}[label=\arabic*.,leftmargin=*]
\item The divisors $S_j$ and their canonical sections give root data
$\OO_X(S_j)^{m_j}=f^*\OO_B(p_j)$, hence the indicated map to the root
curve. Over its order-$m_j$ chart the fibre product is the original
smooth torus family $J'_j$, since its free affine cyclic quotient is
exactly $f^{-1}(D_j)$. The comparison $J'_j\to X$ is finite etale;
in particular pullback is faithful and the stated local comparisons
descend. Away from the two root points the map is the original $f$.
Properness and flatness may be checked after this atlas: properness is
the properness of the original fillings, and flatness holds because a
nonconstant base parameter is a non-zero-divisor in every smooth
total-space local ring. Fibres have dimension two, including the cusp.

\item On a smooth torus chart a vertical field is a holomorphic linear
combination of the translation fields. Its coefficients are constant
on each compact fibre and therefore depend only on the base. The
vertical derivative of an affine transition is its linear part $R_g$;
the translating part, including the logarithmic regluing, has vertical
derivative the identity. At the cusp the two fields
$x_1\partial_{x_1},x_2\partial_{x_2}$ extend as logarithmic vector
fields in every toric chart. The deck maps multiply each $x_i$ by a
base-dependent unit and a power of the base parameter. Their vertical
logarithmic derivative fixes both coefficient directions. The
exponential overlap sends $\partial_{\zeta_i}$ to
$2\pi i x_i\partial_{x_i}$. These transitions depend only on the base
and define $E$; they prove the actual bundle identity
$\mathcal V=\bar f^*E$.

\item We check the kernel and surjectivity at every type of cusp point. At
$t=xy$, the kernel of $df$ is freely generated by
$x\partial_x-y\partial_y$ and $\partial_z$. At $t=xyz$ it is generated
by $x\partial_x-y\partial_y$ and $x\partial_x-z\partial_z$.
Indeed the equation $yz a+xz b+xy c=0$ in the regular local ring forces
$x\mid a,y\mid b,z\mid c$; after division the only relation is that
the three coefficients sum to zero. The logarithmic map sends
$\sum a_i z_i\partial_{z_i}$ to $(\sum a_i)t\partial_t$ and is
surjective. At a finite root point it is the ordinary relative tangent
sequence of the smooth map to $s$, not the singular differential of
$s^{m_j}$. These comparisons prove (D.3), since
$T_{\mathcal B}(-p_0)=\OO(7-12)$.

\item The invariant vector $e_2$ defines a subbundle $\OO\xi\subset E$.
It remains nonzero as a vector of the logarithmic vertical bundle at
the cusp, even where its image in the ordinary tangent fibre is zero.
The quotient line has automorphy factors $-1/\tau$ and $1/\tau$ under
the order-three and order-four generators. The modular discriminant
has weight twelve and no zero in the upper half-plane. Therefore its
reciprocal is a meromorphic section of the twelfth tensor power of
this quotient line. On the cusp chart
$\tau=(\log q)/(2\pi i)+h(q)$ with $h$ holomorphic, so the product
expansion of the discriminant is $q$ times a holomorphic unit.
There are no other zeros or poles, including on the finite root
charts. The quotient line thus has twelfth power
$\OO(-p_0)=\OO(-12)$, and the torsion-free Picard group in \ref{inv:D1.1} makes
it $\OO(-1)$. Its extension by $\OO$ splits because
$H^1(\mathcal B,\OO(1))=0$.

\item For completeness the structure-sheaf computation at the cusp uses the
normalization $P=\mathrm dP_6$, the three double curves and the two
triple points. Their exact normalization complex has global section
complex
\[
 \C\xrightarrow{0}\C^3
 \xrightarrow{(a,b,c)\mapsto(a-b+c,a-b+c)}\C^2.
\]
Exactness at the two local models is the elementary gluing of functions
on the coordinate branches; all higher cohomology of the rational
normalization, the rational double curves and the points vanishes.
The complex therefore gives $(h^0,h^1,h^2)=(1,2,1)$. Its product is
the Cech product of the two-triangle torus and gives a nondegenerate
alternating pairing $H^1\wedge H^1\to H^2$. This also follows directly
from the two opposite oriented triangles: the two coordinate cocycles
have determinant one on their oriented fundamental cycle. These are
the local comparison and product in Claude and Alpoge (2026,
Lemma 9.12), used here on the same normalization and with both triple
points retained.

\item On all other geometric fibres of $\bar f$ the cohomology is that of a
smooth two-torus. Constancy of these three dimensions and proper flat
analytic base change give locally free direct images on the atlas.
Their canonical comparison maps commute with the finite group actions
and with overlap restriction; hence they descend to $\mathcal B$.
The determinant of (D.3), using $W=\bar f^*p_0$, gives
\[
 \omega_{X/\mathcal B}=\bar f^*(\det E)^{-1}.
\]
Relative Serre duality, applied on that same atlas and descended by
its trace pairing, gives $R^2\bar f_*\OO_X=\det E=\OO(-1)$.
Cup product is an isomorphism $\det G\to R^2\bar f_*\OO_X$, since it
is so on every geometric fibre by the preceding calculation.

\item The class $\gamma$ is nowhere zero in $G$. On a smooth marked fibre,
$H^1(\OO)$ is the quotient of constant lattice cocycles by the two
period rows. If $\gamma$ were a combination of those rows, evaluation
on $\widehat w,\widehat\delta$ would make both coefficients zero,
contradicting $\gamma(\widehat\gamma)=1$. The same argument applies
on each finite root cover, and affine translations act trivially on
these classes. At the cusp $\gamma$ vanishes on the angular sublattice
and descends to a nonzero character of its rank-two deck lattice.
In the normalization complex it is the corresponding nonzero
one-cocycle on the two-triangle torus. Its cup product with a complementary
integral character is nonzero, so it remains nonzero there. Thus
$\OO\gamma\subset G$ is a subbundle; its quotient is $\OO(-1)$ by
the determinant just computed. Once again $H^1(\OO(1))=0$ splits the
extension. This proves every assertion in (D.2).
\end{enumerate}
\end{proof}

\begin{proof}[Koszul and direct-image proof of Theorem~\ref{fam:st:A.5.2}]
\begin{enumerate}[label=\arabic*.,leftmargin=*]
\item Put $L=\OO_{\mathcal B}(-5)$. The first connecting map of (D.3) is
the holomorphic bundle map
$\kappa:L\to E\otimes G$. In the splittings fixed in \ref{inv:D1.2} the
coefficient in $\OO\xi\otimes\OO\gamma$ lies in $H^0(\OO(5))=0$.
The remaining coefficients lie in
$H^0(\OO(4))=\C y$, twice, and $H^0(\OO(3))=\C x$.
Therefore, as a map from $E^*$ to $G$ with its $L$ factor understood,
\[
 \kappa=\begin{pmatrix}0&a y\\ b y&c_0x\end{pmatrix}. \tag{D.5}
\]
We use $c_0$ for this coefficient to distinguish it from the deformation
parameter. The logarithmic period derivative at the cusp has leading
term the invertible matrix $B_0/(2\pi i)$. Hence
$\det\kappa=-ab y^2$ is nonzero there, so $ab\ne0$.
At the order-three point $\tau-\rho$ has order one on the root chart:
$j=1728t$ and the elliptic ramification order of $j$ at $\rho$ is
three. The marked torus period derivative is consequently nonzero
there: the two constant identity columns in the normalized period
matrix exclude an infinitesimal linear-frame change that could
otherwise conceal a nonzero derivative. In (D.5), $y=0$ and
$x\ne0$ at that point, so $c_0\ne0$.
These arguments establish the three nonzero coefficients; they do not
infer them from a global cohomology dimension. Rescaling the two
quotient frames and the source frame makes all three coefficients one.

\item The next connecting map is cup product with $\kappa$. In suitable
bases it is
\[
 L\otimes G=\OO(-5)\oplus\OO(-6)
 \xrightarrow{\ \delta_1\ }
 E\otimes\det G=\OO(-1)\oplus\OO(-2),\qquad
 \delta_1=\begin{pmatrix}-y&0\\-x&y\end{pmatrix}.       \tag{D.6}
\]
The sign follows by wedging first with $\gamma$ and then with the
second generator of $G$; reversing both sign conventions leaves the
cokernels unchanged. As its determinant is $-y^2$, it is injective as
a morphism of sheaves. The first connecting map is
$\delta_0=(0,y,y,x)^t$ and is a subbundle inclusion, since $x,y$ have
no common zero. Its cokernel is
\[
 \OO\oplus\OO(-1)\oplus\OO(2).                       \tag{D.7}
\]
Indeed one constant linear combination of the two $\OO(-1)$ summands
is unused, and the remaining quotient is given by the exact sequence
$0\to\OO(-5)\xrightarrow{(y,x)}\OO(-1)\oplus\OO(-2)
\xrightarrow{(x,-y)}\OO(2)\to0$.

\item The cokernel of (D.6) is supported at the order-three root point. There
$x$ is a unit, and elementary row and column operations identify it
with
\[
 Q=\OO(-1)\big|_{2p_1},                               \tag{D.8}
\]
where $2p_1$ means the length-two divisor on the root chart, not the
double coarse point. The generator is the first target component,
and its annihilator is $y^2$. The long direct-image sequence now gives
\[
 R^0\bar f_*\mathcal T=E,\quad
 R^1\bar f_*\mathcal T=\OO\oplus\OO(-1)\oplus\OO(2),
\]
\[
 0\longrightarrow Q\longrightarrow R^2\bar f_*\mathcal T
 \longrightarrow\OO(-6)\longrightarrow0.              \tag{D.9}
\]
This last extension splits, since $\operatorname{Ext}^1(\OO(-6),Q)
=H^1(Q(6))=0$; its splitting is not asserted canonical.

\item Let $\nu:P=\mathrm dP_6\to W\hookrightarrow X$ be the normalization
immersion. There is an exact sequence
\[
 0\longrightarrow\mathcal T\longrightarrow T_X
 \longrightarrow\nu_*K_P\longrightarrow0.             \tag{D.10}
\]
To verify the quotient, in a chart $W=(z_1\cdots z_r=0)$ the map takes
a vector field to its normal coefficient on each of the $r$ separate
branches. The coefficients lift independently; its kernel is precisely
the vector fields tangent to every branch. Thus the quotient is the
pushforward of the normal line of $\nu$, not $\OO_W(W)$.
The form $d\log x_1\wedge d\log x_2\wedge dq$ becomes a unit times
$dz_1\wedge dz_2\wedge dz_3$ in every unimodular chart and is invariant
under every cusp deck transformation. It trivializes $\nu^*K_X$.
Adjunction for the immersion gives its normal line $K_P$, proving
(D.10). Since $P$ is rational, its $K_P$ cohomology is $(0,0,1)$ by
Serre duality. Finite pushforward preserves that cohomology.

\item It remains to identify the extension at $p_0$ in degree two. We give
a fibre-sensitive calculation. The kernel of
\[
 \Omega_X^1|_W\longrightarrow\nu_*\Omega_P^1
\]
is $\nu_*\OO_P\otimes N_W^*$. At $xy=0$ its sections are
$a(y,z)y\,dx+b(x,z)x\,dy$; at $xyz=0$ they are
$a(y,z)yz\,dx+b(x,z)xz\,dy+c(x,y)xy\,dz$.
These identifications glue by the common conormal $dq$.
Since $N_W$ and $K_X|_W$ are trivial and $H^0(P,\Omega_P^1)=0$, the
Serre-dual space of $H^2(W,T_X|_W)$ is one-dimensional.
Top-degree base change identifies this group with the residue fibre
of $(R^2\bar f_*T_X)_{p_0}$. Its generic rank is one by (D.9).
A finite module over the discrete valuation ring $\C\{q\}$ of rank
one and residue dimension one is free of rank one. Thus the extension
by $\C_{p_0}$ increases the degree of its line part by $p_0$, without
adding torsion. We obtain the complete root direct images
\[
 R^0\bar f_*T_X=\OO\oplus\OO(-1),\quad
 R^1\bar f_*T_X=\OO\oplus\OO(-1)\oplus\OO(2),
\]
\[
 R^2\bar f_*T_X=\OO(6)\oplus Q.                       \tag{D.11}
\]
On the order-three chart, an invariant coefficient of the generator
of $Q$ must have exponent $2$ modulo $3$, whereas $Q$ only retains
exponents $0,1$. Therefore $\rho_*Q=0$. Applying the exact $\rho_*$
to (D.11) and using
$d(0)=0,d(-1)=d(2)=-1,d(6)=0$ proves (D.4).
Finally Leray over a curve has only columns zero and one. The usual
cohomology of these displayed line bundles gives the four asserted
dimensions. This is a direct-image computation, not an inference from
the tangent Euler characteristic.
\end{enumerate}
\end{proof}

\begin{proof}[Actual-gluing Kodaira--Spencer calculation for Theorem~\ref{fam:st:A.5.2}]
\begin{enumerate}[label=\arabic*.,leftmargin=*]
\item All line bundles other than the trivial summands in
$R^q\bar f_*\mathcal V=E\otimes R^q\bar f_*\OO_X$ are acyclic:
they have weights $-1,-2$. Hence $H^1(X,\mathcal V)=
\C(\xi\otimes\gamma)$ and its other positive-degree groups vanish.
The three direct images of $\bar f^*\OO(-5)$ have weights
$-5$, $(-5,-6)$, and $-6$, all acyclic by \ref{inv:D1.1}. Consequently (D.3)
induces an isomorphism in every cohomology degree. Sequence (D.10),
whose quotient has no cohomology in degrees zero and one, then gives
the claimed isomorphism $H^1(\mathcal V)\to H^1(T_X)$.

\item We use the Cech convention in which the derivative of a transition
arrow is composed with the inverse of the central arrow. A lattice
transition is translation by $\Pi_c\lambda$, and its derivative in
$c$ is $\gamma(\lambda)\xi$. The linear deck matrices are independent
of $c$. On the finite charts the affine translations contribute
$\gamma(v_j)\xi/m_j$; their $m_j$th iterates give
$\gamma(v_j)\xi$, precisely the lattice cocycle. On the overlap with
the smooth part the derivative of the logarithmic translating section
is $\gamma(v_j)(\log s_j)\xi/(2\pi i)$. Its Cech coboundary accounts
for the difference of these representatives. On the cusp the deck
multiplier differentiates to $\gamma(\lambda)\xi$ in logarithmic
coordinates, and the exponential overlap identifies it with the same
cocycle. Thus the global Kodaira--Spencer class is
$[\xi\otimes\gamma]$, with a minus sign if the tangent-lift convention
$[\bar\partial h]$ is used instead. The sign is recorded; nonvanishing
does not depend on it.

\item Nonvanishing can also be checked on a regular torus. In the quotient
description of $H^1(\OO)$ used above, $[\gamma]\ne0$; the first
coefficient of (D.5) is zero, so $\xi\gamma$ cannot be a scalar
multiple of the base Kodaira--Spencer class. Its image in
$H^1(T_X|_F)$ therefore remains nonzero. Since $H^1(X,T_X)$ is
one-dimensional by \ref{inv:D1.3}, this proves the asserted isomorphism.
\end{enumerate}
\end{proof}

\begin{proof}[Kuranishi ideal calculation for Theorem~\ref{fam:st:A.5.2}]
\begin{enumerate}[label=\arabic*.,leftmargin=*]
\item Let $\alpha$ be the generator in \ref{inv:D2.1}. The constructed holomorphic
family gives an actual convergent deformation tangent to $\alpha$.
Differentiating its integrability equation to second order gives
$\bar\partial\alpha_2=-[\alpha,\alpha]/2$, so the cohomology class
of $[\alpha,\alpha]$ is zero. Bilinearity proves that the bracket
$H^1(T_X)\times H^1(T_X)\to H^2(T_X)$ vanishes on every pair.
This argument alone would not exclude higher obstructions, so we
make the analytic argument as well. A Kuranishi germ is a closed
analytic subgerm of its tangent space $H^1(T_X)=\C$. Its classifying
map from the actual parameter disc has nonzero derivative by \ref{inv:D2.1}.
Composing with this embedding gives a holomorphic map of discs with
nonzero derivative; the inverse function theorem makes its image
contain a neighbourhood of zero in $\C$. Every equation of the
Kuranishi germ vanishes on that neighbourhood and is therefore zero.
The germ, including its analytic structure, is the smooth line.
The same inverse-function argument makes the classifying map an
isomorphism of germs. Pulling back the semiuniversal deformation
therefore realizes every sufficiently small deformation in the
$c$-family. This is semiuniversality; the positive-dimensional
automorphism group is not being discarded or interpreted as a
fine-moduli assertion.
\end{enumerate}
\end{proof}



